%% Question bank chapter file - generated by bank_tools.py build (do not edit while the bank is still being built:
%% it is regenerated from the per-test data in Question Bank\data). Every question is one bankq block:
%%   \begin{bankq}[id=..., source={...}, paper=.., qno=.., type=..., subtopic={...}, answer={...}, answer_src=key|solved,
%%                 working={...}, difficulty=1-10, diff_parts={K S I C F}, diff_note={...},
%%                 group=<passage id or empty>, also_in={duplicate copies in other tests}]
%% difficulty = min(10, K+S+I+C+F-1): K concepts 1-3, S steps 1-4, I insight 0-2, C computation 0-2, F format 0-1
%%   ...question LaTeX, options, figures (\bfig / \bfigi -> img/)...
%%   \end{bankq}
%% Passages shared by several questions are bankpassage blocks; their questions carry group=<passage id>.
\documentclass[10pt]{article}
\usepackage[a4paper,margin=1.8cm]{geometry}
\usepackage[T1]{fontenc}
\usepackage{amsmath,amssymb,mathtools}
\usepackage{graphicx}
\usepackage[export]{adjustbox}
\usepackage{xcolor,needspace,keyval}
\usepackage{enumitem}
\usepackage{booktabs,array,tabularx}
\usepackage[hidelinks]{hyperref}
\setlength{\parindent}{0pt}
\setlength{\parskip}{3pt}
\makeatletter
\define@key{bq}{id}{\def\bq@id{#1}}\define@key{bq}{source}{\def\bq@src{#1}}\define@key{bq}{paper}{\def\bq@paper{#1}}
\define@key{bq}{qno}{\def\bq@qno{#1}}\define@key{bq}{type}{\def\bq@type{#1}}\define@key{bq}{subtopic}{\def\bq@sub{#1}}
\define@key{bq}{answer}{\def\bq@ans{#1}}\define@key{bq}{answer_src}{\def\bq@asrc{#1}}\define@key{bq}{working}{\def\bq@work{#1}}
\define@key{bq}{group}{\def\bq@group{#1}}\define@key{bq}{also_in}{\def\bq@also{#1}}
\define@key{bq}{difficulty}{\def\bq@diff{#1}}\define@key{bq}{diff_parts}{\def\bq@dparts{#1}}\define@key{bq}{diff_note}{\def\bq@dnote{#1}}
\newenvironment{bankq}[1][]{\def\bq@work{}\def\bq@also{}\def\bq@group{}\def\bq@diff{?}\def\bq@dparts{}\def\bq@dnote{}\setkeys{bq}{#1}%
  \par\needspace{4\baselineskip}\medskip\noindent{\small\textbf{\bq@id}\hfill\color{gray}[\bq@src, Paper \bq@paper, Q\bq@qno]
  \bq@type\ $\cdot$ \bq@sub\ $\cdot$ \textbf{D\bq@diff}}\par\nopagebreak\relax}
  {\par{\small\color{teal}Answer: \textbf{\bq@ans}\ifx\bq@work\@empty\else\ (solved: \bq@work)\fi
  \ifx\bq@also\@empty\else\par\color{gray}Also in: \bq@also\fi}\par}
\define@key{bp}{id}{\def\bp@id{#1}}\define@key{bp}{source}{\def\bp@src{#1}}\define@key{bp}{questions}{\def\bp@qs{#1}}
\newenvironment{bankpassage}[1][]{\setkeys{bp}{#1}\par\needspace{4\baselineskip}\bigskip\noindent
  {\small\textbf{Passage \bp@id}\hfill\color{gray}[\bp@src; questions \bp@qs]}\par\nopagebreak\relax\itshape}{\par}
\makeatother
\newcommand{\bfig}[2]{\par\smallskip{\centering\includegraphics[width=#1pt,max width=0.62\linewidth,max height=0.32\textheight]{img/#2}\par}\smallskip}
\newcommand{\bfigi}[2]{\includegraphics[width=#1pt,max width=\linewidth,max height=0.32\textheight,valign=c]{img/#2}}
\newcommand{\bmissing}[1]{\textit{[figure missing in the source conversion: \detokenize{#1}]}}
\newenvironment{qoptions}{\begin{itemize}[leftmargin=1.8em,itemsep=1pt,topsep=2pt,label={}]}{\end{itemize}}
\newcommand{\opt}[2]{\item[(#1)] #2}
\providecommand{\fwstart}{}\providecommand{\fwend}{}\providecommand{\qfullsplit}{}
\newcolumntype{L}[1]{>{\raggedright\arraybackslash}p{#1}}
\newcolumntype{C}[1]{>{\centering\arraybackslash}p{#1}}
\newcolumntype{R}[1]{>{\raggedleft\arraybackslash}p{#1}}
\providecommand{\sectionhdr}[1]{}\providecommand{\subsectionhdr}[1]{}


\begin{document}
{\Large\bfseries Chemical Kinetics --- Question Bank}\par\smallskip
{\small Chemistry chapter C09. Questions from all converted test papers, grouped by source series. Answers: printed key where it matches the paper, otherwise solved (answer\_src=solved, with working). Exact repeats across papers appear once, with also\_in listing the other copies.}\par\medskip\hrule

\section*{Score-I papers}
\begin{bankq}[id=SI01-P1-CHEM-S2-Q12, source={Score-I - Test 01 (16 Jun 2021)}, paper=1, qno=12, type=multi, subtopic={Arrhenius equation}, answer={A,B,C,D}, answer_src=key, difficulty=8, diff_parts={K3 S1 I2 C2 F1}, diff_note={concepts: standard problem in this chapter, arrhenius equation, apply the relevant results | steps: apply the relevant results; compute}]
Which of the following statements are correct ?
\begin{qoptions}
\opt{A}{$\dfrac{k_p}{k_a}=e^{\Delta E/RT}$ where $k_p$ = Rate constant for reaction with catalyst.

$K_a$ = Rate constant for reaction without catalyst

$\Delta E$ = Difference in activation energy for the two ways of reactions.}
\opt{B}{Alcohol acts as a negative catalyst for the conversion of $\mathrm{Na_2SO_3}\longrightarrow \mathrm{Na_2SO_4}$ .}
\opt{C}{TEL (tetraethyl lead) acts as a negative catalyst in the knocking of petrol.}
\opt{D}{Dilute $\mathrm{H_2SO_4}$ acts as a positive catalyst in the hydrolysis of esters to carboxylic acids and alcohols.}
\end{qoptions}
\end{bankq}
\begin{bankq}[id=SI01-P1-CHEM-S3-Q05, source={Score-I - Test 01 (16 Jun 2021)}, paper=1, qno=5, type=numerical, subtopic={Rate laws}, answer={7}, answer_src=key, difficulty=8, diff_parts={K3 S1 I2 C2 F1}, diff_note={concepts: standard problem in this chapter, rate laws, apply the relevant results | steps: apply the relevant results; compute}]
An optically active substance `A' is decomposing into optically active substance `B' and `C' as
\par\noindent $\mathrm{A}\xrightarrow{K=10^{-3}\,\mathrm{min}^{-1}}2\mathrm{B}+\mathrm{C}$\par
The specific rotations of A, B and C are $+40^\circ$, $+10^\circ$ and $-30^\circ$ per mole, respectively. Initially `A' and `C' were present in $4:3$ mole ratio. Close to how many hours the sample becomes optically inactive ?

[Given : $\ell\mathrm{n}\,\dfrac{7}{5}=0.34,\ \ell\mathrm{n}\,\dfrac{20}{13}=0.42$ ]
\end{bankq}
\begin{bankq}[id=SI01-P2-CHEM-S2-Q11, source={Score-I - Test 01 (16 Jun 2021)}, paper=2, qno=11, type=multi, subtopic={Rate laws}, answer={A,B,C}, answer_src=key, difficulty=8, diff_parts={K3 S1 I2 C2 F1}, diff_note={concepts: standard problem in this chapter, rate laws, apply the relevant results | steps: apply the relevant results; compute}]
Select incorrect statements :
\begin{qoptions}
\opt{A}{Pre--exponential factor for zero order reaction is a unitless quantity.}
\opt{B}{If $t_{1/4}=30$ sec then $t_{1/2}=60$ sec for first order reaction}
\opt{C}{If $t_{1/3}=30$ sec then $t_{2/3}=90$ sec for sec--order reaction}
\opt{D}{If $t_{1/5}=30$ sec then $t_{3/5}=90$ sec for zero order reaction.}
\end{qoptions}
\end{bankq}
\begin{bankq}[id=SI07-P1-CHEM-S1-Q06, source={Score-I - Test 07 (04 Feb 2022)}, paper=1, qno=6, type=multi, subtopic={Rate laws}, answer={A,B,D}, answer_src=key, difficulty=8, diff_parts={K3 S1 I2 C2 F1}, diff_note={concepts: standard problem in this chapter, rate laws, apply the relevant results | steps: apply the relevant results; compute}]
A reaction $\mathrm{A_2+B_2\rightleftharpoons2AB}$ follows the path :

$\mathrm{A_2}\ \overset{\text{fast}}{\rightleftharpoons}\ \mathrm{A+A}$

$\mathrm{A+B_2}\ \xrightarrow{\text{slow}}\ \mathrm{AB+B}$

$\mathrm{A+B}\ \overset{\text{fast}}{\rightleftharpoons}\ \mathrm{AB}$
\begin{qoptions}
\opt{A}{overall order of the reaction is 1.5}
\opt{B}{order w.r.t. A$_2$ is 0.5}
\opt{C}{order w.r.t. B$_2$ is 0.5}
\opt{D}{Order of a reaction may change by changing conditions like concentration, temperature or catalyst.}
\end{qoptions}
\end{bankq}
\begin{bankq}[id=SI07-P1-CHEM-S2-Q02, source={Score-I - Test 07 (04 Feb 2022)}, paper=1, qno=2, type=numerical, subtopic={Order}, answer={1}, answer_src=key, difficulty=8, diff_parts={K3 S1 I2 C2 F1}, diff_note={concepts: standard problem in this chapter, order, apply the relevant results | steps: apply the relevant results; compute}]
For parallel reaction of reactant X to products Y \& Z given :
\bfig{52.5}{SI07_p1_chem_s2_q02_1.png}

The ratio $\dfrac{\mathrm k_1}{\mathrm k_2}=\dfrac23$ [k = rate constant]

When all of X is completely consumed then find mole \% of Z obtained?
\end{bankq}
\begin{bankq}[id=SI07-P1-CHEM-S3-Q06, source={Score-I - Test 07 (04 Feb 2022)}, paper=1, qno=6, type=numerical, subtopic={Arrhenius equation}, answer={9.00}, answer_src=key, difficulty=8, diff_parts={K3 S1 I2 C2 F1}, diff_note={concepts: standard problem in this chapter, arrhenius equation, apply the relevant results | steps: apply the relevant results; compute}]
A reaction follows temperature - rate constant relation as

$\ell\mathrm n\,\mathrm k=30-\dfrac{5000}{\mathrm T}+3\,\ell\mathrm n(3\mathrm T)$

If activation energy = 11.8 kcal/mol and R = 2 cal k$^{-1}$ mol$^{-1}$ then find the value of $\dfrac{\mathrm T}{100}$ where T = reaction temperature [in kelvin].
\end{bankq}
\begin{bankq}[id=SI09-P2-CHEM-S1-Q03, source={Score-I - Test 09 (16 Feb 2022)}, paper=2, qno=3, type=multi, subtopic={Rate laws}, answer={B,C,D}, answer_src=key, difficulty=8, diff_parts={K3 S1 I2 C2 F1}, diff_note={concepts: standard problem in this chapter, rate laws, apply the relevant results | steps: apply the relevant results; compute}]
For the following reaction\\
$2\mathbf X+\mathbf Y\xrightarrow{\;k\;}\mathbf P$\\
the rate of reaction is $\dfrac{d[\mathrm P]}{dt}=k[\mathbf X]$. Two moles of $\mathbf X$ are mixed with one mole of $\mathbf Y$ to make 1.0 L of solution. At 50 s, 0.5 mole of $\mathbf Y$ is left in the reaction mixture. The correct statement(s) about the reaction is(are)\\
(Use: $\ln 2=0.693$)
\begin{qoptions}
\opt{A}{The rate constant, $k$, of the reaction is $13.86\times10^{-4}\ \mathrm{s^{-1}}$.}
\opt{B}{Half--life of $\mathbf X$ is 50 s.}
\opt{C}{At 50 s, $-\dfrac{d[\mathbf X]}{dt}=13.86\times10^{-3}\ \mathrm{mol\,L^{-1}\,s^{-1}}$.}
\opt{D}{At 100 s, $-\dfrac{d[\mathbf Y]}{dt}=3.46\times10^{-3}\ \mathrm{mol\,L^{-1}\,s^{-1}}$}
\end{qoptions}
\end{bankq}
\begin{bankq}[id=SI12-P1-CHEM-S1-Q07, source={Score-I - Test 12 (09 Mar 2022)}, paper=1, qno=7, type=multi, subtopic={Nuclear chemistry}, answer={A,B,C}, answer_src=key, difficulty=3, diff_parts={K1 S2 I0 C0 F1}, diff_note={concepts: nuclear chemistry statements | steps: nuclear reaction balance; temperature independence of decay constant; isodiaphers; average life = 1.44 t1/2}]
Select the correct statement(s)
\begin{qoptions}
\opt{A}{bombarding ${}_7$N$^{14}$ nuclei with $\alpha$-particle, the nuclei of the product formed after release of proton would be ${}_8$O$^{17}$}
\opt{B}{Decay constant does not depend upon temperature}
\opt{C}{Nuclide and it's decay product after $\alpha$-emission are called isodiaphers}
\opt{D}{Half-life of radium is 1580 years. It's average life will be 1097.22 years.}
\end{qoptions}
\end{bankq}

\section*{Score-II papers}
\begin{bankpassage}[id=SII01-P1-CHEM-S2-G7-8, source={Score-II - Test 01 (21 Mar 2021)}, questions={7-8}]
Dependence of rate constant of a reaction is given by Arrhenius equation $k = A\cdot e^{-E_a/RT}$ but originally it's derived from its differential equation.
\end{bankpassage}
\begin{bankq}[id=SII01-P1-CHEM-S2-Q07, source={Score-II - Test 01 (21 Mar 2021)}, paper=1, qno=7, type=single, subtopic={Rate laws}, answer={B}, answer_src=key, difficulty=4, diff_parts={K2 S1 I1 C1 F0}, diff_note={concepts: standard problem in this chapter, rate laws | steps: apply the relevant results; compute}, group=SII01-P1-CHEM-S2-G7-8]
Reactions A + B $\to$ C + D follows following rate law: rate $= k[A]^{+1/2}[B]^{1/2}$. Starting with initial conc. of one mole of A and B each, what is the time taken for amount of A to become 0.25 mole. Given $k = 2.31\times10^{-3}$ sec$^{-1}$.
\begin{qoptions}
\opt{A}{300 sec.}
\opt{B}{600 sec.}
\opt{C}{900 sec.}
\opt{D}{None of these}
\end{qoptions}
\end{bankq}
\begin{bankq}[id=SII01-P1-CHEM-S2-Q08, source={Score-II - Test 01 (21 Mar 2021)}, paper=1, qno=8, type=single, subtopic={Rate laws}, answer={C}, answer_src=key, difficulty=4, diff_parts={K2 S1 I1 C1 F0}, diff_note={concepts: standard problem in this chapter, arrhenius equation | steps: apply the relevant results; compute}, group=SII01-P1-CHEM-S2-G7-8]
For a reaction, the dependence of rate constant (sec$^{-1}$) on temperature (K) is given as
\[\ln K = 10 - \frac{2500}{T} + 3\ln T\]
The activation energy of reaction is
\begin{qoptions}
\opt{A}{$2500+3T$}
\opt{B}{$2500\,R$}
\opt{C}{$(2500+3T)R$}
\opt{D}{$(2500-3\ln T)R$}
\end{qoptions}
\end{bankq}
\begin{bankq}[id=SII02-P1-CHEM-S1-Q04, source={Score-II - Test 02 (28 Mar 2021)}, paper=1, qno=4, type=single, subtopic={Arrhenius equation}, answer={B}, answer_src=key, difficulty=4, diff_parts={K2 S1 I1 C1 F0}, diff_note={concepts: standard problem in this chapter, arrhenius equation | steps: apply the relevant results; compute}]
\textbf{Statement-I:} On increasing the temperature of an elementary biomolecular gaseous reaction by $10^\circ$C occuring at constant volume, the rate of reaction increase substantially by factor of 2 to 3.

\textbf{Statement-2:} The increase in temperature increases the collision frequency (total number of biomolecular collisions per unit time per unit volume) for the above condition.

\begin{qoptions}
\opt{A}{Statement-1 is true, Statement-2 is true and statement-2 is correct explanation for statement-1.}
\opt{B}{Statement-1 is true, Statement-2 is true and statement-2 is NOT correct explanation for statement-1.}
\opt{C}{Statement-1 is true, statement-2 is false.}
\opt{D}{Statement-1 is false, statement-2 is true.}
\end{qoptions}
\end{bankq}
\begin{bankq}[id=SII02-P1-CHEM-S2-Q07, source={Score-II - Test 02 (28 Mar 2021)}, paper=1, qno=7, type=multi, subtopic={Order}, answer={B,D}, answer_src=key, difficulty=8, diff_parts={K3 S1 I2 C2 F1}, diff_note={concepts: standard problem in this chapter, order, apply the relevant results | steps: apply the relevant results; compute}]
Which of the following statements is/are incorrect w.r.t enzyme catalysed reactions.

\begin{qoptions}
\opt{A}{The cavities on enzyme ``Maltase'' \& structure of molecules of Maltose have complimentary shapes for conversion of Maltase to glucose.}
\opt{B}{Just like other catalysts, catalytic action of enzyme is more at a lower temperature.}
\opt{C}{In presence of vitamins \& metal ions; enzyme activity generally increases.}
\opt{D}{Most of the enzyme catalysed reactions are elementary.}
\end{qoptions}
\end{bankq}
\begin{bankq}[id=SII04-P1-CHEM-S3-Q03, source={Score-II - Test 04 (07 May 2021)}, paper=1, qno=3, type=numerical, subtopic={Order}, answer={0.25}, answer_src=key, difficulty=8, diff_parts={K3 S1 I2 C2 F1}, diff_note={concepts: standard problem in this chapter, order, apply the relevant results | steps: apply the relevant results; compute}]
\begin{center}
\bfigi{172.3}{SII04_p1chem_q3_graph.png}
\end{center}

If three reactions:
\[A\longrightarrow\text{Products}\qquad\text{Rate}=K_1[A]^X\]
\[B\longrightarrow\text{Products}\qquad\text{Rate}=K_2[B]^Y\]
\[C\longrightarrow\text{Products}\qquad\text{Rate}=K_3[C]^Z\]
Value of $\left(\dfrac{X+Y}{4Z}\right)$ is?
\end{bankq}
\begin{bankq}[id=SII06-P1-CHEM-S1-Q03, source={Score-II - Test 06 (05 Jul 2022)}, paper=1, qno=3, type=single, subtopic={Order}, answer={C}, answer_src=key, difficulty=4, diff_parts={K2 S1 I1 C1 F0}, diff_note={concepts: standard problem in this chapter, order | steps: apply the relevant results; compute}]
The oxidation of certain metal is found to obey the equation $A^2 = \alpha t + \beta$, where A is the thickness of metal oxide film at time t, $\alpha$ and $\beta$ are constants. The order of this reaction is
\begin{qoptions}
\opt{A}{0}
\opt{B}{1}
\opt{C}{$-1$}
\opt{D}{2}
\end{qoptions}
\end{bankq}
\begin{bankq}[id=SII07-P1-CHEM-S2-Q10, source={Score-II - Test 07 (10 Jul 2022)}, paper=1, qno=10, type=numerical, subtopic={Rate laws}, answer={3}, answer_src=key, difficulty=8, diff_parts={K3 S1 I2 C2 F1}, diff_note={concepts: standard problem in this chapter, rate laws, apply the relevant results | steps: apply the relevant results; compute}]
For a reaction\\
$\mathrm{A(g)}\xrightarrow{\text{Rate constant}=k}\mathrm{P(g)}+\mathrm{Q(g)}$\\
carried in closed and rigid container at constant temperature $T$, following informations are provided.
\begin{center}
\begin{tabular}{|l|c|c|c|}\hline
 & Initial Concentration & $T_{1/2}$ & $T_{3/4}$\\\hline
A(g) & $1\times10^{-2}$ M & 10 Min & 30 Min\\\hline
\end{tabular}
\end{center}
If initial rate of reaction in mole/lit/min is represented in scientific notation as $x\times10^{-y}$ then find the value of $y$.
\end{bankq}
\begin{bankq}[id=SII08-P1-CHEM-S2-Q07, source={Score-II - Test 08 (05 Feb 2023)}, paper=1, qno=7, type=single, subtopic={Rate laws}, answer={A}, answer_src=key, difficulty=4, diff_parts={K2 S1 I1 C1 F0}, diff_note={concepts: standard problem in this chapter, rate laws | steps: apply the relevant results; compute}]
Match the rate expressions in LIST-I for the decomposition of X with the corresponding profiles provided in LIST-II. X$_s$ and k constants having appropriate units.

\begin{center}
\begin{tabular}{|c|p{5.2cm}|c|l|}
\hline
\multicolumn{2}{|c|}{LIST-I} & \multicolumn{2}{c|}{LIST-II} \\
\hline
(I) & rate $= \dfrac{k[X]}{X_s+[X]}$ under all possible initial concentration of X & (P) & \bfigi{85.2}{SII08_p1_chem_s2_q07_1.png} \\
\hline
(II) & rate $= \dfrac{k[X]}{X_s+[X]}$ where initial concentration of X are much less than $X_s$ & (Q) & \bfigi{88.1}{SII08_p1_chem_s2_q07_2.png} \\
\hline
(III) & rate $= \dfrac{k[X]}{X_s+[X]}$ where initial concentration of X are much higher than $X_s$ & (R) & \bfigi{86.6}{SII08_p1_chem_s2_q07_3.png} \\
\hline
(IV) & rate $= \dfrac{k[X]^2}{X_s+[X]}$ where initial concentration of X is much higher than $X_s$ & (S) & \bfigi{81.5}{SII08_p1_chem_s2_q07_4.png} \\
\hline
& & (T) & \bfigi{83.1}{SII08_p1_chem_s2_q07_5.png} \\
\hline
\end{tabular}
\end{center}

\begin{qoptions}
\opt{A}{I$\to$R; II$\to$Q; III$\to$S; IV$\to$T}
\opt{B}{I$\to$R; II$\to$S; III$\to$S; IV$\to$T}
\opt{C}{I$\to$P; II$\to$Q; III$\to$Q; IV$\to$R}
\opt{D}{I$\to$P; II$\to$S; III$\to$Q; IV$\to$R}
\end{qoptions}
\end{bankq}
\begin{bankq}[id=SII12-P1-CHEM-S4-Q03, source={Score-II - Test 12 (04 Feb 2024)}, paper=1, qno=3, type=numerical, subtopic={Arrhenius equation}, answer={1}, answer_src=key, difficulty=8, diff_parts={K3 S1 I2 C2 F1}, diff_note={concepts: standard problem in this chapter, arrhenius equation, apply the relevant results | steps: apply the relevant results; compute}]
The plot of $\log k_f$ versus $1/T$ for a reversible reaction $A(g)\rightleftharpoons P(g)$ is shown.
\bfig{127.6}{SII12_p1_chem_s4_q03_1.png}
Pre-exponential factors for the forward and backward reactions are $10^{15}\ s^{-1}$ and $10^{11}\ s^{-1}$, respectively. If the value of log K for the reaction at 500 K is 5, the value of $|\log k_b|$ at 250 K is \_\_\_\_\_\_. [K = equilibrium constant of the reaction, $k_f$ = rate constant of forward reaction, $k_b$ = rate constant of backward reaction]
\end{bankq}
\begin{bankq}[id=SII13-P1-CHEM-S2-Q01, source={Score-II - Test 13 (06 Feb 2024)}, paper=1, qno=1, type=numerical, subtopic={Rate laws}, answer={2}, answer_src=key, difficulty=8, diff_parts={K3 S1 I2 C2 F1}, diff_note={concepts: standard problem in this chapter, rate laws, apply the relevant results | steps: apply the relevant results; compute}]
For a first order reaction, time required for 99.0\% completion is `x' times for the time required for the completion of 90\% of the reaction then value of x is ?
\end{bankq}
\begin{bankq}[id=SII14-P1-CHEM-S2-Q05, source={Score-II - Test 14 (06 Feb 2024)}, paper=1, qno=5, type=multi, subtopic={Rate laws}, answer={B,C}, answer_src=key, difficulty=8, diff_parts={K3 S1 I2 C2 F1}, diff_note={concepts: standard problem in this chapter, rate laws, apply the relevant results | steps: apply the relevant results; compute}]
For the reaction\\
$Br^-(aq)+BrO_3^-(aq)+H^+(aq)\longrightarrow Br_2(aq)+H_2O(\ell)$\\
the rate expression was found to be $-d[BrO_3^-]/dt=k[Br^-][H^+]^2[BrO_3^-]$\\
Which of the following statement(s) is/are correct ?
\begin{qoptions}
\opt{A}{Molecularity of reaction is 4.}
\opt{B}{Dimensions of pseudo rate constant of the reaction in a buffer solution is $M^{-1}min^{-1}$.}
\opt{C}{Doubling the concentration of all the reactants at the same time will increase the reaction rate by a factor of 16.}
\opt{D}{Rate of decrease of $BrO_3^-$ and rate of formation of $Br_2$ are the same.}
\end{qoptions}
\end{bankq}
\begin{bankq}[id=SII17-P1-CHEM-S1-Q02, source={Score-II - Test 17 (13 Feb 2024)}, paper=1, qno=2, type=single, subtopic={Order}, answer={A}, answer_src=solved, working={First order: k1=2.303 log10/t=2.303/t; k2=2.303 log100/(2t)=2.303/t; ratio 1.}, difficulty=4, diff_parts={K2 S1 I1 C1 F0}, diff_note={concepts: standard problem in this chapter, order | steps: apply the relevant results; compute}]
In the following first order reactions\\
(A) $\xrightarrow{K_1}$ Product, (B) $\xrightarrow{K_2}$ Product , find the ratio $K_1/K_2$ if 90\% of (A) has been reacted in time `t' while 99\% of (B) has been reacted in time 2t :-
\begin{qoptions}
\opt{A}{1}
\opt{B}{2}
\opt{C}{1/2}
\opt{D}{$2\left(\dfrac{1-2\log3}{2-\log11-2\log3}\right)$}
\end{qoptions}
\end{bankq}
\begin{bankq}[id=SII17-P1-CHEM-S2-Q01, source={Score-II - Test 17 (13 Feb 2024)}, paper=1, qno=1, type=numerical, subtopic={Arrhenius equation}, answer={140}, answer_src=solved, working={k=(k1k2/k3)\^{}(2/3): Ea=(2/3)(180+80-50)=140 kJ/mol.}, difficulty=8, diff_parts={K3 S1 I2 C2 F1}, diff_note={concepts: standard problem in this chapter, arrhenius equation, apply the relevant results | steps: apply the relevant results; compute}]
For a complex reaction A $\xrightarrow{k}$ products\\
$E_{a_1}$=180 kJ/mol; $E_{a_2}$ =80 kJ/mol; $E_{a_3}$ = 50 kJ/mol\\
Overall rate constant k is related to individual rate constant by the equation $k=\left(\dfrac{k_1.k_2}{k_3}\right)^{2/3}$\\
Activation energy (kJ/mol) for the overall reaction is :-
\end{bankq}
\begin{bankq}[id=SII19-P1-CHEM-S1-Q09, source={Score-II - Test 19 (20 Feb 2024)}, paper=1, qno=9, type=single, subtopic={Order}, answer={B}, answer_src=key, difficulty=4, diff_parts={K2 S1 I1 C1 F0}, diff_note={concepts: standard problem in this chapter, order | steps: apply the relevant results; compute}]
Half life of a first order reaction :\\
A(aq.) $\rightarrow$ 2B(aq.) is 40 min at 27$^\circ$C and 20 min at 37$^\circ$C. The activation energy of reaction is ($\ln 2 = 0.7$, $R = 2.0$ cal mol$^{-1}$K$^{-1}$)
\begin{qoptions}
\opt{A}{19530 cal/ mol}
\opt{B}{13020 cal/ mol}
\opt{C}{26040 cal/ mol}
\opt{D}{6510 cal/ mol}
\end{qoptions}
\end{bankq}
\begin{bankq}[id=SII19-P1-CHEM-S2-Q04, source={Score-II - Test 19 (20 Feb 2024)}, paper=1, qno=4, type=numerical, subtopic={Order}, answer={12}, answer_src=key, difficulty=8, diff_parts={K3 S1 I2 C2 F1}, diff_note={concepts: standard problem in this chapter, order, apply the relevant results | steps: apply the relevant results; compute}]
The half-life of a first order reaction is one hour. In what time will the concentration be reduced to $\dfrac78$ of its initial value :- ($\log 7 = 0.84$)
\end{bankq}
\begin{bankq}[id=SII21-P1-CHEM-S2-Q07, source={Score-II - Test 21 (14 Mar 2024)}, paper=1, qno=7, type=single, subtopic={Rate laws}, answer={D}, answer_src=key, difficulty=4, diff_parts={K2 S1 I1 C1 F0}, diff_note={concepts: standard problem in this chapter, rate laws | steps: apply the relevant results; compute}]
Match the rate expressions in LIST-I for the decomposition of X with the corresponding profiles provided in LIST-II. X$_s$ and k constants having appropriate units.

\begin{center}
\begin{tabular}{|c|p{5.5cm}|c|c|}
\hline
\multicolumn{2}{|c|}{LIST-I} & \multicolumn{2}{c|}{LIST-II} \\
\hline
(I) & rate $= \dfrac{k[X]}{X_s+[X]}$ under all possible initial concentration of X & (P) & \bfigi{88.6}{SII21_p1_chem_s2_q07_1.png} \\
\hline
(II) & rate $= \dfrac{k[X]}{X_s+[X]}$ where initial concentration of X are much less than $X_s$ & (Q) & \bfigi{91.9}{SII21_p1_chem_s2_q07_2.png} \\
\hline
(III) & rate $= \dfrac{k[X]}{X_s+[X]}$ where initial concentration of X are much higher than $X_s$ & (R) & \bfigi{90.5}{SII21_p1_chem_s2_q07_3.png} \\
\hline
(IV) & rate $= \dfrac{k[X]^2}{X_s+[X]}$ where initial concentration of X is much higher than $X_s$ & (S) & \bfigi{84.9}{SII21_p1_chem_s2_q07_4.png} \\
\hline
& & (T) & \bfigi{85.0}{SII21_p1_chem_s2_q07_5.png} \\
\hline
\end{tabular}
\end{center}

\begin{qoptions}
\opt{A}{I$\to$R; II$\to$S; III$\to$S; IV$\to$T}
\opt{B}{I$\to$P; II$\to$Q; III$\to$Q; IV$\to$R}
\opt{C}{I$\to$R; II$\to$S; III$\to$Q; IV$\to$R}
\opt{D}{I$\to$R; II$\to$Q; III$\to$S; IV$\to$T}
\end{qoptions}
\end{bankq}
\begin{bankq}[id=SII22-P1-CHEM-S1-Q02, source={Score-II - Test 22 (14 Mar 2024)}, paper=1, qno=2, type=single, subtopic={Rate laws}, answer={D}, answer_src=key, difficulty=4, diff_parts={K2 S1 I1 C1 F0}, diff_note={concepts: standard problem in this chapter, rate laws | steps: apply the relevant results; compute}]
Which of the following is true statement ?
\begin{qoptions}
\opt{A}{Rate of reaction does not changes on changing stoichiometry of reaction}
\opt{B}{Rate can be negative}
\opt{C}{If rate of appearance of $NH_3$ is 1.7 kg/hr in the reaction $N_2+H_2\longrightarrow NH_3$ then rate of disappearance of $H_2$ will be 42.5 gm/min}
\opt{D}{Rate of disappearance of reactants or rate of appearance of products does not changes on changing stoichiometry of the reaction}
\end{qoptions}
\end{bankq}
\begin{bankq}[id=SII23-P1-CHEM-S1-Q01, source={Score-II - Test 23 (17 Mar 2024)}, paper=1, qno=1, type=multi, subtopic={Order}, answer={A,D}, answer_src=key, difficulty=8, diff_parts={K3 S1 I2 C2 F1}, diff_note={concepts: standard problem in this chapter, order, apply the relevant results | steps: apply the relevant results; compute}]
SO$_3$ gas is entering the environment at a constant rate of $6.93\times10^{-6}$ g/L/day due to emission of polluting gases from thermal power plant at Kota but at the same time it is decomposing and following first-order kinetics with half-life of 100 days. Based on these details, select the correct statement(s) from the following.
\begin{qoptions}
\opt{A}{Concentration of SO$_3$ in Kota is $1.25\times10^{-5}$ M (assume SO$_3$ present in air reaches steady state).}
\opt{B}{If $10^3$ L of air is passed through 1 L pure water (assuming all SO$_3$ is dissolved in it) and the resulting solution is titrated against 1 N NaOH solution, 15 ml of NaOH solution is required to reach the end point.}
\opt{C}{An industry is manufacturing H$_2$SO$_4$ at the rate of 980 kg per day with the use of SO$_3$ in air and it should use $8\times10^5$ L air per day.}
\opt{D}{If SO$_3$ emission is stopped, then after 1000 days, its concentration will reduce to $1.22\times10^{-8}$ M.}
\end{qoptions}
\end{bankq}
\begin{bankq}[id=SII23-P1-CHEM-S1-Q02, source={Score-II - Test 23 (17 Mar 2024)}, paper=1, qno=2, type=multi, subtopic={Rate laws}, answer={C,D}, answer_src=key, difficulty=8, diff_parts={K3 S1 I2 C2 F1}, diff_note={concepts: standard problem in this chapter, rate laws, apply the relevant results | steps: apply the relevant results; compute}]
Decomposition of A$_2$(g) and B$_3$(g) follows first-order kinetics as follows.
\[A_2(g)\to 2A(g);\ K_1(h^{-1}) = e^{-\frac{14000(J)}{RT}+5}\]
\[B_3(g)\to 3B(g);\ K_2(h^{-1}) = e^{-\frac{20000(J)}{RT}+10}\]
Here, K$_1$ and K$_2$ are the rate constants with respect to disappearance of A$_2$ and B$_3$, respectively. One mole, each of A$_2$(g) and B$_3$(g), is taken in a 100 L evacuated flask and at some temperature at which they start decomposing at the same rate. The incorrect information regarding the reactions is
\begin{qoptions}
\opt{A}{The temperature at which the reactions are performed is $\dfrac{1200}{8.314}$K.}
\opt{B}{At any instant, $\left(\dfrac{P_{A_2}}{P_{B_3}}\right)$ will be constant and equal to 1.0, where P$_{A_2}$ and P$_{B_3}$ are the partial pressures of gases at that time.}
\opt{C}{At some instant, the total pressure of all gases in the flask may be less than 0.2 atm.}
\opt{D}{At any instant, $\left(\dfrac{P_A}{P_B}\right)$ will be constant and equal to 1.0, where P$_A$ and P$_B$ are the partial pressures of gases at that time.}
\end{qoptions}
\end{bankq}
\begin{bankq}[id=SII24-P1-CHEM-S2-Q05, source={Score-II - Test 24 (17 Mar 2024)}, paper=1, qno=5, type=numerical, subtopic={Order}, answer={400}, answer_src=key, difficulty=8, diff_parts={K3 S1 I2 C2 F1}, diff_note={concepts: standard problem in this chapter, order, apply the relevant results | steps: apply the relevant results; compute}]
For a particular reaction A $\to$ product, graph of reciprocal of concentration of reactant (y-axis) versus time (x-axis) is a straight line having intercept 10 (on y-axis) and slope 0.1. Calculate time (in min) required for 80\% completion of reaction. [Time - min, concentration = molar]
\end{bankq}
\begin{bankq}[id=SII25-P1-CHEM-S1-Q16, source={Score-II - Test 25 (20 Mar 2024)}, paper=1, qno=16, type=single, subtopic={Rate laws}, answer={C}, answer_src=key, difficulty=4, diff_parts={K2 S1 I1 C1 F0}, diff_note={concepts: standard problem in this chapter, rate laws | steps: apply the relevant results; compute}]
For a reaction $2P+Q\rightarrow$ product, rate law is :\\
$-\dfrac{d[P]}{dt}=k[P]$\\
At a time when $t=\dfrac2k$, concentration of the reactant is `P' : ($C_0$ = initial concentration of reactant `P')
\begin{qoptions}
\opt{A}{$\dfrac{C_0}{e}$}
\opt{B}{$eC_0$}
\opt{C}{$\dfrac{C_0}{e^2}$}
\opt{D}{$\dfrac1{C_0}$}
\end{qoptions}
\end{bankq}
\begin{bankpassage}[id=SII26-P1-CHEM-S3-G5-6, source={Score-II - Test 26 (20 Mar 2024)}, questions={5-6}]
The mixing of different gases by frequent colllision is called diffusion and the process of escaping gas molecules through tiny orifice in vaccum is called effusion.
\end{bankpassage}
\begin{bankq}[id=SII26-P1-CHEM-S3-Q05, source={Score-II - Test 26 (20 Mar 2024)}, paper=1, qno=5, type=numerical, subtopic={Rate laws}, answer={17.65}, answer_src=key, difficulty=8, diff_parts={K3 S1 I2 C2 F1}, diff_note={concepts: standard problem in this chapter, rate laws, apply the relevant results | steps: apply the relevant results; compute}, group=SII26-P1-CHEM-S3-G5-6]
The reaction P$_{(g)}\to2$Q$_{(g)}+$R$_{(g)}$ follows I$^{\text{st}}$ order kinetics having rate constant 0.023 sec$^{-1}$. 1 mole compound P is taken in an evacuted sealed flask. After `t' time in the flask a hole is generated and the gases effuses according to Graham's law. The initially effusing gaseous mixture contains double the moles of Q than P. Determine value of t (in sec) if gram molecular masses of P, Q and R are 24, 6, 9 respectively.\\
\relax[Given : log3/2 = 0.176, ln 10 = 2.3]
\end{bankq}
\begin{bankq}[id=SII26-P1-CHEM-S3-Q06, source={Score-II - Test 26 (20 Mar 2024)}, paper=1, qno=6, type=numerical, subtopic={Rate laws}, answer={17.00}, answer_src=key, difficulty=8, diff_parts={K3 S1 I2 C2 F1}, diff_note={concepts: standard problem in this chapter, rate laws, apply the relevant results | steps: apply the relevant results; compute}, group=SII26-P1-CHEM-S3-G5-6]
The mass \% of R in the mixture coming out initially ? (Give nearest integer as your answer).
\end{bankq}
\begin{bankq}[id=SII28-P1-CHEM-S3-Q02, source={Score-II - Test 28 (23 Mar 2024)}, paper=1, qno=2, type=numerical, subtopic={Rate laws}, answer={1380}, answer_src=key, difficulty=8, diff_parts={K3 S1 I2 C2 F1}, diff_note={concepts: standard problem in this chapter, rate laws, apply the relevant results | steps: apply the relevant results; compute}]
For reaction A(g) + B(g) $\rightleftharpoons$ C(g), at 300 K, the rate of backward reaction at time, t, is 10 times of forward reaction for the reaction. Calculate free energy change (in cal) involved at that time of reaction.\\
(Given : lnx = 2.3 logx)
\end{bankq}
\begin{bankq}[id=SII29-P1-CHEM-S2-Q06, source={Score-II - Test 29 (26 Mar 2024)}, paper=1, qno=6, type=multi, subtopic={Order}, answer={B,D}, answer_src=key, difficulty=8, diff_parts={K3 S1 I2 C2 F1}, diff_note={concepts: standard problem in this chapter, order, apply the relevant results | steps: apply the relevant results; compute}]
For the reaction,
A $\to$ Products (order = n)
the time of completion of reaction is

\begin{qoptions}
\opt{A}{Infinite for all value of `n'}
\opt{B}{Infinite for n = 1}
\opt{C}{$\dfrac{[A_0]^{1-n}}{K(1-n)}$ for n $\neq$ 1}
\opt{D}{$\dfrac{[A_0]^{1-n}}{K(1-n)}$ for n < 1}
\end{qoptions}
\end{bankq}
\begin{bankq}[id=SII30-P1-CHEM-S1-Q03, source={Score-II - Test 30 (27 Mar 2024)}, paper=1, qno=3, type=single, subtopic={Order}, answer={A}, answer_src=key, difficulty=4, diff_parts={K2 S1 I1 C1 F0}, diff_note={concepts: standard problem in this chapter, order | steps: apply the relevant results; compute}]
The half-life of radioisotopic bromine - 82 is 36 hours. The fraction which remains after one day is \underline{\hspace{2em}} $\times10^{-2}$.\\
(Given antilog 0.2006 = 1.587)
\begin{qoptions}
\opt{A}{63}
\opt{B}{126}
\opt{C}{32}
\opt{D}{100}
\end{qoptions}
\end{bankq}
\begin{bankq}[id=SII31-P1-CHEM-S1-Q17, source={Score-II - Test 31 (29 Mar 2024)}, paper=1, qno=17, type=single, subtopic={Order}, answer={C}, answer_src=key, difficulty=4, diff_parts={K2 S1 I1 C1 F0}, diff_note={concepts: standard problem in this chapter, order | steps: apply the relevant results; compute}]
At 30$^\circ$C, the half-life for the decomposition of AB$_2$ is 200 sec and is independent of the initial concentration of AB$_2$. The time required (in sec) for 80\% of the AB$_2$ to decompose is (Given: log 2 = 0.30; log 3 = 0.48)
\begin{qoptions}
\opt{A}{200 sec}
\opt{B}{323 sec}
\opt{C}{467 sec}
\opt{D}{532 sec}
\end{qoptions}
\end{bankq}
\begin{bankq}[id=SII32-P1-CHEM-S1-Q04, source={Score-II - Test 32 (29 Mar 2024)}, paper=1, qno=4, type=single, subtopic={Rate laws}, answer={B}, answer_src=key, difficulty=4, diff_parts={K2 S1 I1 C1 F0}, diff_note={concepts: standard problem in this chapter, rate laws | steps: apply the relevant results; compute}]
In the following reaction
$$x_1A + y_1B \to x_2C + y_2D$$
$$\resizebox{\linewidth}{!}{$\log_{10}\left[-\dfrac{dA}{dt}\right]+0.301=\log_{10}\left[-\dfrac{dB}{dt}\right]+0.477=\log_{10}\left[\dfrac{dC}{dt}\right]+0.602=\log_{10}\left[\dfrac{dD}{dt}\right]+0.699$}$$
Then find $x_1,y_1,x_2$ and $y_2$ respectively

\begin{qoptions}
\opt{A}{2, 3, 4 and 5}
\opt{B}{30, 20, 15, 12}
\opt{C}{3, 2, 5 and 4}
\opt{D}{30, 20, 25, 24}
\end{qoptions}
\end{bankq}
\begin{bankq}[id=SII34-P1-CHEM-S3-Q08, source={Score-II - Test 34 (02 Feb 2025)}, paper=1, qno=8, type=numerical, subtopic={Arrhenius equation}, answer={5.75}, answer_src=key, difficulty=8, diff_parts={K3 S1 I2 C2 F1}, diff_note={concepts: standard problem in this chapter, arrhenius equation, apply the relevant results | steps: apply the relevant results; compute}]
In a 1st order reaction, rate constant varies with temperature according to the graph given below.
\bfig{116.5}{SII34_p1_chem_s3_q08_1.png}
At 27$^\circ$C, $1.0\times10^{-4}$ \% of the reactant molecules are able to cross the potential barrier. At 72$^\circ$C, the slope of this graph is equal to 0.2 K$^{-1}$ sec$^{-1}$. Calculate the value of rate constant (in s$^{-1}$) at 72$^\circ$C, assuming that activation energy does not change in this temperature range.\\
\relax[Given : R = $\dfrac{25}{3}$ J/K - mol, ln10 = 2.3]
\end{bankq}
\begin{bankq}[id=SII40-P1-CHEM-S4-Q03, source={Score-II - Test 40 (16 Mar 2025)}, paper=1, qno=3, type=numerical, subtopic={Arrhenius equation}, answer={6}, answer_src=key, difficulty=8, diff_parts={K3 S1 I2 C2 F1}, diff_note={concepts: standard problem in this chapter, arrhenius equation, apply the relevant results | steps: apply the relevant results; compute}]
The plot of log $k_f$ versus $\dfrac1T$ for a reversible reaction A (g) $\rightleftharpoons$ P (g) is shown.
\bfig{138.7}{SII40_p1_chem_s4_q03_1.png}
Pre-exponential factors for the forward and backward reactions are $10^{15}$ s$^{-1}$ and $10^{11}$ s$^{-1}$, respectively. If the value of log $K_b$ for the reaction at 250 K is $-5$, the value of log K at 500 K is \underline{\hspace{1cm}}.\\
\relax[K = equilibrium constant of the reaction\\
$k_f$ = rate constant of forward reaction\\
$k_b$ = rate constant of backward reaction]
\end{bankq}
\begin{bankq}[id=SII41-P1-CHEM-S1-Q18, source={Score-II - Test 41 (19 Mar 2025)}, paper=1, qno=18, type=single, subtopic={Rate laws}, answer={C}, answer_src=key, difficulty=4, diff_parts={K2 S1 I1 C1 F0}, diff_note={concepts: standard problem in this chapter, rate laws | steps: apply the relevant results; compute}]
What will be the expression of rate constant ($k$) of reaction in terms of pressure for following reaction, when the initial pressure is $P_0$ and, at time `$t$', total pressure is $P_T$ (considering reaction is following first order kinetics) :\\
\centerline{2A(g) $\longrightarrow$ B(g)}
\begin{qoptions}
\opt{A}{$k=\dfrac1t\,\ell\text{n}\dfrac{P_0}{2P_0-P_T}$}
\opt{B}{$k=\dfrac1{2t}\,\ell\text{n}\dfrac{P_0}{P_0-P_T}$}
\opt{C}{$k=\dfrac1{2t}\,\ell\text{n}\dfrac{P_0}{2P_T-P_0}$}
\opt{D}{$k=\dfrac{2.303}t\,\ell\text{og}\dfrac{P_0}{P_T}$}
\end{qoptions}
\end{bankq}
\begin{bankq}[id=SII43-P1-CHEM-S1-Q07, source={Score-II - Test 43 (23 Mar 2025)}, paper=1, qno=7, type=single, subtopic={Rate laws}, answer={A}, answer_src=key, difficulty=4, diff_parts={K2 S1 I1 C1 F0}, diff_note={concepts: standard problem in this chapter, rate laws | steps: apply the relevant results; compute}]
For a gaseous reaction :\\
CH$_3$CHO(g) $\xrightarrow{\ \Delta\ }$ CO(g) + CH$_4$(g)\\
initial pressure is 80 mm of Hg and total pressure at the end of 20 minutes is 120 mm of Hg. The rate constant of the reaction assuming first order kinetics is - (ln2 = 0.693)
\begin{qoptions}
\opt{A}{$3.465\times10^{-2}$ min$^{-1}$}
\opt{B}{34.65 min$^{-1}$}
\opt{C}{3.465 min$^{-1}$}
\opt{D}{0.3465 min$^{-1}$}
\end{qoptions}
\end{bankq}
\begin{bankq}[id=SII45-P1-CHEM-S1-Q14, source={Score-II - Test 45 (26 Mar 2025)}, paper=1, qno=14, type=single, subtopic={Order}, answer={B}, answer_src=key, difficulty=4, diff_parts={K2 S1 I1 C1 F0}, diff_note={concepts: standard problem in this chapter, order | steps: apply the relevant results; compute}]
\textbf{Statement-1 :} The hydrolysis of ethyl acetate is a second order reaction, in reality.\\
CH$_3$COOC$_2$H$_5$ + H$_2$O $\to$ CH$_3$COOH + C$_2$H$_5$OH\\
$r=K[\mathrm{CH_3COOC_2H_5}]\,[\mathrm{H_2O}]$\\
\textbf{Statement-2 :} As the concentration of water does not get altered much during the course of reaction, the reaction behaves as first order reaction with respect to ester.
\begin{qoptions}
\opt{A}{Statement-1 and Statement 2 are correct and Statement 2 explains Statement 1}
\opt{B}{Statement-1 and Statement-2 are correct and Statement -2 doesn't explain Statement 1}
\opt{C}{Statement-1 is correct and Statement-2 is incorrect}
\opt{D}{Statement-1 and Statement-2 both are incorrect.}
\end{qoptions}
\end{bankq}
\begin{bankq}[id=SII45-P1-CHEM-S2-Q03, source={Score-II - Test 45 (26 Mar 2025)}, paper=1, qno=3, type=numerical, subtopic={Order}, answer={2}, answer_src=key, difficulty=8, diff_parts={K3 S1 I2 C2 F1}, diff_note={concepts: standard problem in this chapter, order, apply the relevant results | steps: apply the relevant results; compute}]
A first order reaction completes 75\% in 40 minutes. In what time (in hours), the concentration of reactant will reduce from 0.64 M to 0.01 M ?
\end{bankq}
\begin{bankq}[id=SII46-P1-CHEM-S1-Q08, source={Score-II - Test 46 (26 Mar 2025)}, paper=1, qno=8, type=multi, subtopic={Rate laws}, answer={A,D}, answer_src=key, difficulty=8, diff_parts={K3 S1 I2 C2 F1}, diff_note={concepts: standard problem in this chapter, rate laws, apply the relevant results | steps: apply the relevant results; compute}]
$\alpha$ - maltose (C$_{12}$H$_{22}$O$_{11}$) can be hydrolysed to glucose (C$_6$H$_{12}$O$_6$) according to the following reaction:
\[\mathrm{C_{12}H_{22}O_{11}(aq) + H_2O(\ell) \rightarrow 2\,C_6H_{12}O_6(aq)}\]
Given: Standard enthalpy of formation of C$_{12}$H$_{22}$O$_{11}$(aq) = $-2238$ kJ/mol; Standard enthalpy of formation of H$_2$O($\ell$) = $-285$ kJ/mol; Standard enthalpy of formation of C$_6$H$_{12}$O$_6$(aq) = $-1263$ kJ/mol
\begin{center}
\begin{tabular}{|l|c|c|c|}
\hline
Time (min) & 0 & 50 & 100 \\
\hline
Conc.\ of $\alpha$-maltose (M) & 4.0 & 1.0 & 0.25 \\
\hline
\end{tabular}
\end{center}
which of the following statement(s) is/are true?
\begin{qoptions}
\opt{A}{The hydrolysis of $\alpha$-maltose is exothermic.}
\opt{B}{Heat liberated in combustion of 1 mole of $\alpha$-maltose(s) must be greater than the heat liberated in combustion of 2 moles of glucose(s).}
\opt{C}{Increasing temperature will increase the degree of hydrolysis of $\alpha$-maltose.}
\opt{D}{The hydrolysis of $\alpha$-maltose follows $1^{\mathrm{st}}$ order kinetics.}
\end{qoptions}
\end{bankq}
\begin{bankq}[id=SII47-P1-CHEM-S2-Q02, source={Score-II - Test 47 (30 Mar 2025)}, paper=1, qno=2, type=numerical, subtopic={Order}, answer={3}, answer_src=key, difficulty=8, diff_parts={K3 S1 I2 C2 F1}, diff_note={concepts: standard problem in this chapter, order, apply the relevant results | steps: apply the relevant results; compute}]
For a second order reaction, containing single reactant : A $\to$ B, it is given that :\\
$T_{75\%}=X\cdot t_{50\%}$\\
Find the value of X
\end{bankq}

\section*{Score Advanced papers}
\begin{bankq}[id=SA02-P1-CHEM-S3-Q01, source={Score Advanced - Test 02 (13 May 2021)}, paper=1, qno=1, type=numerical, subtopic={Arrhenius equation}, answer={3}, answer_src=solved, working={k doubles from 300 to 400 K (ln k2/k1=0.7), so t1/2=1.5 min; 2 to 0.5 bar is two half lives}, difficulty=7, diff_parts={K2 S2 I2 C1 F1}, diff_note={concepts: first order reaction at two temperatures, arrhenius equation | steps: rate constant ratio from activation energy; half life; time}]
For the first order reaction A(g) $\to$ 2B(g) + C(g) the half life for decomposition of A is 3 minute at 300K. Calculate the time (in minute) in which partial pressure of A(g) will drop from 2 bar to 0.5 bar at 400K. Given activation energy of reaction is 840 R. (Take ln2 = 0.7)
\end{bankq}
\begin{bankq}[id=SA02-P2-CHEM-S3-Q03, source={Score Advanced - Test 02 (13 May 2021)}, paper=2, qno=3, type=numerical, subtopic={Arrhenius equation}, answer={2}, answer_src=solved, working={Integrate Ea/RT\^{}2 from 200 to 400 K: (1/2)(2+(20/7) ln2)=2}, difficulty=7, diff_parts={K2 S2 I2 C1 F1}, diff_note={concepts: rate constant ratio with variable activation energy, arrhenius equation | steps: Arrhenius derivative; integrate; evaluate}]
The activation energy of a reaction is given as E$_a$ (cal mol$^{-1}$) = 0.01T$^2$ + $\dfrac{20}{7}$T where T is temperature in kelvin. Find value of $\ell$n$\dfrac{k_2}{k_1}$ where k$_1$ and k$_2$ are rate constants at temperature 200 K and 400 K respectively. (Use 2 = 0.7 and R = 2 cal mol$^{-1}$ k$^{-1}$)
\end{bankq}
\begin{bankq}[id=SA05-P1-CHEM-S1-Q16, source={Score Advanced - Test 05 (15 Jul 2021)}, paper=1, qno=16, type=single, subtopic={Rate laws}, answer={C}, answer_src=key, difficulty=5, diff_parts={K2 S1 I2 C1 F0}, diff_note={concepts: ester hydrolysis by titration, rate laws | steps: acid formed is V\_t-V\_0; total possible V\_inf-V\_0}]
For acid catalysed hydrolysis of ester, rate constant is given as K= $\dfrac{1}{t}\ln\left(\dfrac{V_\infty - V_0}{V_\infty - V_t}\right)$ where V$_0$, V$_t$ and V$_\infty$ are volume of standard NaOH solution required to neutralise acid present at given time. If ester gets 50\% hydrolysed then -
\begin{qoptions}
\opt{A}{V$_\infty$ = V$_t$}
\opt{B}{V$_\infty$ = V$_t$--V$_0$}
\opt{C}{V$_\infty$ = 2V$_t$--V$_0$}
\opt{D}{V$_\infty$ = V$_t$ --2V$_0$}
\end{qoptions}
\end{bankq}
\begin{bankq}[id=SA06-P1-CHEM-S3-Q02, source={Score Advanced - Test 06 (04 Sep 2021)}, paper=1, qno=2, type=numerical, subtopic={Rate laws}, answer={5}, answer_src=key, difficulty=5, diff_parts={K2 S1 I1 C1 F1}, diff_note={concepts: statements on kinetics, rate laws | steps: check each of six statements}]
How many of the following statements regarding kinetics of a chemical reaction are correct.\par
(i) The probability of collision of three or more molecules with proper orientation at same time is very rare thats why reactions having molecularity three or more are very rare.\par
(ii) In a second order reaction the number of reactant can be one but the rate will be proportional to the square of its concentration.\par
(iii) The rate of a chemical reaction at any time is always proportional to the initial concentration of the reactant molecules raised to the power equal to their order.\par
(iv) The rate of a chemical reaction also depends upon the strength of bond of reactant molecules and the number of bond to be broken.\par
(v) The total kinetic energy of the colliding molecules must be greater than some minimum value.\par
(vi) The colliding particles must be properly oriented in space when they collide.
\end{bankq}
\begin{bankq}[id=SA08-P1-CHEM-S1-Q06, source={Score Advanced - Test 08 (22 Apr 2023)}, paper=1, qno=6, type=single, subtopic={Rate laws}, answer={B}, answer_src=key, difficulty=5, diff_parts={K2 S1 I2 C1 F0}, diff_note={concepts: order from a table of rates, rate laws | steps: compare experiments with one or two changes}]
For the following reaction : 2A + B + C $\longrightarrow$ 3D, rate was measured at different concentrations and noted. Identify the option which is correct on the basis of data.\par
{\scriptsize\setlength{\tabcolsep}{1pt}\noindent\begin{tabular}{|c|c|c|c|c|}
\hline
Exp. & conc.of [A] & conc. of [B] & conc. of [C] & Rate of formation of [D] \\ \hline
(1) & 0.03 M & 0.04 M & 0.02 M & 10$^{-3}$ M min$^{-1}$ \\ \hline
(2) & 0.06 M & 0.08 M & 0.02 M & 10$^{-3}$ M min$^{-1}$ \\ \hline
(3) & 0.06 M & 0.16 M & 0.04 M & 5 $\times$ 10$^{-4}$ M min$^{-1}$ \\ \hline
(4) & 0.06 M & 0.08 M & 0.04 M & 10$^{-3}$ M min$^{-1}$ \\ \hline
\end{tabular}}\par
\begin{qoptions}
\opt{A}{Overall order of the reaction is negative}
\opt{B}{Order of the reaction with respect to A will be 1.}
\opt{C}{Overall order of the reaction is 2}
\opt{D}{Order of the reaction with respect to B will be 1.}
\end{qoptions}
\end{bankq}
\begin{bankq}[id=SA08-P2-CHEM-S1-Q06, source={Score Advanced - Test 08 (22 Apr 2023)}, paper=2, qno=6, type=multi, subtopic={Arrhenius equation}, answer={A,B,C}, answer_src=key, difficulty=8, diff_parts={K3 S1 I2 C2 F1}, diff_note={concepts: temperature coefficient of the rate constant, arrhenius equation, sigma=Ea/RT\^{}2 | steps: sigma=Ea/RT\^{}2; solve for Ea at given T}]
In order to study effect of temperature on rate of reaction, another term ``$\sigma$'' is defined as $\sigma = \dfrac{1}{\text{K}}\dfrac{\text{dk}}{\text{dT}}$ [where k is rate constant and T represents temperature]. If for a reaction value of $\sigma = \dfrac{1.25\times10^6}{\text{T}^3}\ \text{K}^{-1}$ [T should be in Kelvin], then identify the correct statement(s).
\begin{qoptions}
\opt{A}{Activation energy for the given reaction at 250 Kelvin is 10 Kcal/mol.}
\opt{B}{Activation energy for the given reaction at 2000 Kelvin is 1.25 Kcal/mol.}
\opt{C}{Rate of increase of rate constant with temperature is higher at lower temperature range than at higher temperature range.}
\opt{D}{The value of $\dfrac{\text{dK}}{\text{dT}}$ at 2000 Kelvin is 0.625 $\times$ K.}
\end{qoptions}
\end{bankq}
\begin{bankq}[id=SA09-P2-CHEM-S3-Q03, source={Score Advanced - Test 09 (29 Apr 2023)}, paper=2, qno=3, type=numerical, subtopic={Order}, answer={120}, answer_src=key, difficulty=8, diff_parts={K3 S1 I2 C2 F1}, diff_note={concepts: second order rate constant and time for conversion, t1/2=1/(k a0), time for 60 percent | steps: t1/2=1/(k a0); time for 60 percent}]
The decomposition of A$_2$B$_2$ follows the rate equation r = K[A$_2$B$_2$]$^2$.\par
A$_2$B$_2$ $\xrightarrow{\text{K}}$ A$_2$ + B$_2$. The kinetic data of the reaction is given below.\par
{\footnotesize\setlength{\tabcolsep}{2pt}\noindent\begin{tabular}{|L{5cm}|c|}
\hline
Initial conc. of A$_2$B$_2$ (mole/lit) & 7\\ \hline
First half life of A$_2$B$_2$ in sec. & 1440\\ \hline
\end{tabular}}\par
If the initial concentration of A$_2$B$_2$ were 56 m mol/lit, Calculate the time (in sec) taken for the conversion of A$_2$B$_2$ to reach 60\% of initial value.\par
(Report your answer as $\dfrac{\text{t}}{1000}$ )
\end{bankq}
\begin{bankq}[id=SA10-P2-CHEM-S3-Q06, source={Score Advanced - Test 10 (06 May 2023)}, paper=2, qno=6, type=numerical, subtopic={Order}, answer={3}, answer_src=key, difficulty=8, diff_parts={K3 S1 I2 C2 F1}, diff_note={concepts: order with respect to a reactant from pressure data, compare the rates at different pressures | steps: compare the rates at different pressures}]
For the reaction : 2NO(g) + H$_2$(g) $\rightarrow$ N$_2$O(g) + H$_2$O(g), the value of $\dfrac{-\text{dp}}{\text{dt}}$ was found to be 1.5 PaS$^{-1}$ for a pressure of 372 Pa for NO and 0.25 PaS$^{-1}$ for a pressure of 152 Pa for NO, the pressure of H$_2$ being constant. If pressure of NO was kept costant, the value of $\dfrac{-\text{dp}}{\text{dt}}$ was found to be 1.60 PaS$^{-1}$ for a pressure of 289 Pa for H$_2$ and 0.79PaS$^{-1}$ for a pressure of 144 Pa for H$_2$. If order of reaction with respect to NO and H$_2$ are a and b, respectively, then find nearest integer value of (a + b) ?
\end{bankq}
\begin{bankq}[id=SA11-P1-CHEM-S3-Q08, source={Score Advanced - Test 11 (09 May 2023)}, paper=1, qno=8, type=numerical, subtopic={Order}, answer={40}, answer_src=key, difficulty=8, diff_parts={K3 S1 I2 C2 F1}, diff_note={concepts: time of completion from a graph, order, read the order from the graph | steps: read the order from the graph; time from the rate law}, also_in={Score Advanced - Test 29 (01 May 2024) P1 Q8}]
In a chemical reaction where A(g) is converting to B(g), the following observation is made. What is the time of completion for the reaction (in min)?\par
\bfig{146.7}{SA11_p1_chem_s3_q08_1.png}
\end{bankq}
\begin{bankq}[id=SA14-P1-CHEM-S1-Q06, source={Score Advanced - Test 14 (16 May 2023)}, paper=1, qno=6, type=multi, subtopic={Rate laws}, answer={B,C,D}, answer_src=key, difficulty=8, diff_parts={K3 S1 I2 C2 F1}, diff_note={concepts: rate law of 2X + Y -> P, rate laws, order from the data and the rate expressions | steps: order from the data and the rate expressions}, also_in={Score Advanced - Test 20 (06 Apr 2024) P1 Q6; TOAS - Test 04 (15 Mar 2022) P1 Q8}]
For the following reaction 2X + Y $\xrightarrow{\text{k}}$ P . The rate of reaction is $\dfrac{\text{d [P]}}{\text{dt}}$ = k [X].\par
Two moles of X are mixed with one mole of Y to make 1.0 L of solution. At 50 s, 0.5 mole of Y is left in the reaction mixture. The correct statement(s) about the reaction is/are (\textbf{Use}: ln2 = 0.693)
\begin{qoptions}
\opt{A}{The rate constant, k, of the reaction is 13.86 $\times$ 10$^{-4}$ s$^{-1}$}
\opt{B}{Half-life of X is 50 s}
\opt{C}{At 50s, $-\dfrac{\text{d [X]}}{\text{dt}}$ = 13.86 $\times$ 10$^{-3}$ mol L$^{-1}$ s$^{-1}$}
\opt{D}{At, 100s, $-\dfrac{\text{d [Y ]}}{\text{dt}}$ = 3.465 $\times$ 10$^{-3}$ mol L$^{-1}$ s$^{-1}$}
\end{qoptions}
\end{bankq}
\begin{bankq}[id=SA14-P1-CHEM-S2-Q09, source={Score Advanced - Test 14 (16 May 2023)}, paper=1, qno=9, type=single, subtopic={Rate laws}, answer={A}, answer_src=key, difficulty=7, diff_parts={K3 S1 I2 C2 F0}, diff_note={concepts: fractional lives of reactions of various orders, rate laws, t(x/y) relations for first and other orders | steps: t(x/y) relations for first and other orders}]
For the reaction: A $\rightarrow$ product, t$_{\text{x/y}}$ represents the times in which x/y fraction of reactant is converted into products.\par
{\footnotesize\setlength{\tabcolsep}{1.5pt}\noindent\begin{tabular}{|c|c|c|L{5.3cm}|}
\hline
 & List-I & & List-II \\
\hline
(P) & t$_{5/9}$ & (1) & Equal to 54s, if t$_{1/3}$ is 18s in case of first-order reaction \\
\hline
(Q) & t$_{19/27}$ & (2) & Equal to 32s, if t$_{1/4}$ is 16s in case of first-order reaction \\
\hline
(R) & t$_{7/8}$ & (3) & Equal to 56s, if t$_{1/3}$ is 4s in case of second-order reaction \\
\hline
(S) & t$_{7/16}$ & (4) & Equal to 30s, if t$_{1/3}$ is 18s in case of zero-order reaction. \\
\hline
 & & (5) & Equal to 28s, if t$_{1/2}$ is 16s in case of zero-order reaction. \\
\hline
\end{tabular}}\par
Codes:
\begin{qoptions}
\opt{A}{P $\rightarrow$ 4;Q $\rightarrow$ 1;R $\rightarrow$ 5;S $\rightarrow$ 2}
\opt{B}{P $\rightarrow$ 4;Q $\rightarrow$ 1;R $\rightarrow$ 2;S $\rightarrow$ 5}
\opt{C}{P $\rightarrow$ 4;Q $\rightarrow$ 2;R $\rightarrow$ 5;S $\rightarrow$ 3}
\opt{D}{P $\rightarrow$ 1;Q $\rightarrow$ 4;R $\rightarrow$ 2;S $\rightarrow$ 5}
\end{qoptions}
\end{bankq}
\begin{bankq}[id=SA16-P1-CHEM-S1-Q01, source={Score Advanced - Test 16 (20 May 2023)}, paper=1, qno=1, type=single, subtopic={Rate laws}, answer={B}, answer_src=key, difficulty=8, diff_parts={K3 S2 I2 C2 F0}, diff_note={concepts: rate constant of hydrogen peroxide decomposition by titration, rate laws, volume strength to molarity | steps: volume strength to molarity; permanganate titration; first order law}]
20 V H$_2$O$_2$ solution is kept in the open vessel and is allowed to decompose by first order kinetics. After 6 hours, 11.35 ml of this H$_2$O$_2$ sample was taken and diluted to 100 ml by adding water. 10 ml of this diluted solution requires 25 ml of 0.025 M KMnO$_4$ in acidic medium for titration. Calculate the rate constant for the decomposition of H$_2$O$_2$. [Given : ln 2 = 0.70, ln 5 = 1.60]
\begin{qoptions}
\opt{A}{0.04 hr$^{-1}$}
\opt{B}{0.05 hr$^{-1}$}
\opt{C}{0.15 hr$^{-1}$}
\opt{D}{0.30 hr$^{-1}$}
\end{qoptions}
\end{bankq}
\begin{bankq}[id=SA17-P1-CHEM-S2-Q09, source={Score Advanced - Test 17 (25 May 2023)}, paper=1, qno=9, type=single, subtopic={Rate laws}, answer={A}, answer_src=key, difficulty=8, diff_parts={K3 S2 I2 C2 F0}, diff_note={concepts: matching kinetics quantities, rate laws, stoichiometric ratios | steps: stoichiometric ratios; half life order; first order time; Nernst}]
Match the list\par
{\footnotesize\setlength{\tabcolsep}{2pt}\noindent\begin{tabular}{|c|L{4.6cm}|c|c|}\hline
\multicolumn{2}{|c|}{\textbf{List-I}} & \multicolumn{2}{c|}{\textbf{List-II}}\\\hline
(P) & Ratio of rate of appearance (gm/min) of SO$_3$(g) and rate of disappearance (gm/min) of SO$_2$(g) for : 2SO$_2$(g) + O$_2$(g) $\longrightarrow$ 2SO$_3$(g) ; would be & (1) & 0\\\hline
(Q) & Based on following the observations for a reaction A $\xrightarrow{k}$ product\par [A]$_0$ (in M) : 0.0677\quad 0.136\quad 0.272\par t$_{1/2}$ (in min) : 2\quad 4\quad 8\par the order of reaction would be & (2) & 3\\\hline
(R) & Time (min) taken by a reaction to decrease the concentration of the reactant(R) from 0.2 mole/litre to 0.025 mole/litre if its rate constant corresponding to rate of disappearance of R is 0.693 min$^{-1}$, would be & (3) & 1.25\\\hline
(S) & If emf at 25\textdegree C of the following Galvanic cell :\par Pt $|$H$_2$(g)$|$ CH$_3$COOH(aq.)$_{(C_1)}$ $\|$ HCOOH(aq.)$_{(C_2)}$ $|$ H$_2$(g)$|$ Pt\par is 0.0177 Volt, then ratio of [H$^+$] in anodic compartment and cathodic compartment would be approximately & (4) & 2\\\hline
\end{tabular}}\par
\begin{qoptions}
\opt{A}{P $\rightarrow$ 3;Q $\rightarrow$ 1;R $\rightarrow$ 2;S $\rightarrow$ 4}
\opt{B}{P $\rightarrow$ 2;Q $\rightarrow$ 2;R $\rightarrow$ 3;S $\rightarrow$ 1}
\opt{C}{P $\rightarrow$ 1;Q $\rightarrow$ 3;R $\rightarrow$ 3;S $\rightarrow$ 4}
\opt{D}{P $\rightarrow$ 3;Q $\rightarrow$ 2;R $\rightarrow$ 1;S $\rightarrow$ 4}
\end{qoptions}
\end{bankq}
\begin{bankq}[id=SA18-P1-CHEM-S1-Q01, source={Score Advanced - Test 18 (27 May 2023)}, paper=1, qno=1, type=multi, subtopic={Rate laws}, answer={B,D}, answer_src=key, difficulty=8, diff_parts={K3 S1 I2 C2 F1}, diff_note={concepts: rate of change of pressure of reactants and products, rate laws, rates relative to stoichiometry | steps: rates relative to stoichiometry; total pressure}]
For the reaction at 273$^\circ$C\par
N$_2$(g) + 3H$_2$(g) $\longrightarrow$ 2NH$_3$(g)\par
It is observed that the pressure of H$_2$(g) falls from 760 mmHg to 380 mmHg in 100 sec.\par
Select the correct statement(s).
\begin{qoptions}
\opt{A}{Average rate of reaction is $5\times10^{-3}$ atm/sec.}
\opt{B}{Average rate of formation of NH$_3$(g) is $\frac{1}{300}$ atm/sec.}
\opt{C}{Average rate of disappearance of N$_2$(g) is twice of rate of appearance of NH$_3$(g).}
\opt{D}{$\log\left(\frac{-\Delta P_{H_2}}{\Delta t}\right)=\log\left(\frac{-\Delta P_{N_2}}{\Delta t}\right)+0.4771$}
\end{qoptions}
\end{bankq}
\begin{bankq}[id=SA18-P2-CHEM-S3-Q01, source={Score Advanced - Test 18 (27 May 2023)}, paper=2, qno=1, type=numerical, subtopic={Arrhenius equation}, answer={1}, answer_src=key, difficulty=5, diff_parts={K2 S1 I1 C1 F1}, diff_note={concepts: activation energy from ln k against 1/T, arrhenius equation | steps: slope=-Ea/R}]
A reaction follows temperature - rate constant relation as\par
$\ell nk=30-\frac{500}{T}$\par
If $R=2$ cal k$^{-1}$ mol$^{-1}$ then the value of activation energy of reaction (in kcal/mole) is :
\end{bankq}
\begin{bankq}[id=SA18-P2-CHEM-S3-Q02, source={Score Advanced - Test 18 (27 May 2023)}, paper=2, qno=2, type=numerical, subtopic={Arrhenius equation}, answer={6}, answer_src=key, difficulty=8, diff_parts={K3 S1 I2 C2 F1}, diff_note={concepts: activation energy with a catalyst, arrhenius equation, equal rate constants at the two temperatures | steps: equal rate constants at the two temperatures}]
A first order chemical reaction.A $\rightarrow$ P is occurring at 500 K. If the same reaction is to be carried out in presence of a catalyst, temperature of 400 K is required to be maintained in order to maintain the same rate of reaction. If the catalyst lowers the activation energy by 30 kJ, then the value of activation energy of the original pathway (uncatalysed reaction) is $x$ kJ/mol, then $\frac{x}{25}$ is
\end{bankq}
\begin{bankq}[id=SA19-P1-CHEM-S3-Q04, source={Score Advanced - Test 19 (04 Apr 2024)}, paper=1, qno=4, type=numerical, subtopic={Order}, answer={2}, answer_src=key, difficulty=8, diff_parts={K3 S1 I2 C2 F1}, diff_note={concepts: half life from iodometric titration of SO2, order, moles of SO2 from the titration | steps: moles of SO2 from the titration; first order rate law}]
0.01 moles of sulphuryl chloride SO$_2$Cl$_2$(g) is taken in a sealed tube \& heated to 400 K.\par
It decomposes by following 1st order kinetics according to reaction:\par
SO$_2$Cl$_2$(g) $\rightarrow$ SO$_2$(g) + Cl$_2$(g)\par
The tube is broken after 4 hours \& gas passed through a 25 ml of an acidified 1 N iodine solution where all SO$_2$ is oxidized to SO$_4^{2-}$ (consider other components are not involved in reaction with iodine). The resulting solution required 10 ml of 1 M Hypo solution. What is half life in hours.
\end{bankq}
\begin{bankq}[id=SA20-P1-CHEM-S2-Q09, source={Score Advanced - Test 20 (06 Apr 2024)}, paper=1, qno=9, type=single, subtopic={Rate laws}, answer={A}, answer_src=key, difficulty=7, diff_parts={K3 S1 I2 C2 F0}, diff_note={concepts: fractional life times for a reaction, rate laws, integrated rate law for each order | steps: integrated rate law for each order}]
For the reaction: A $\rightarrow$ product, $t_{x/y}$ represents the times in which $x/y$ fraction of reactant is converted into products.\par
{\footnotesize\setlength{\tabcolsep}{2pt}\noindent\begin{tabular}{|c|L{1.0cm}|c|L{5.2cm}|}
\hline
 & \multicolumn{1}{c|}{List-I} & & \multicolumn{1}{c|}{List-II}\\ \hline
(P) & $t_{5/9}$ & (1) & Equal to 54s, if $t_{1/3}$ is 18s in case of first-order reaction\\ \hline
(Q) & $t_{19/27}$ & (2) & Equal to 32s, if $t_{1/4}$ is 16s in case of first-order reaction\\ \hline
(R) & $t_{7/8}$ & (3) & Equal to 56s, if $t_{1/3}$ is 4s in case of second-order reaction\\ \hline
(S) & $t_{7/16}$ & (4) & Equal to 30s, if $t_{1/3}$ is 18s in case of zero-order reaction.\\ \hline
 & & (5) & Equal to 28s, if $t_{1/2}$ is 16s in case of zero-order reaction.\\ \hline
\end{tabular}}\par
Codes:
\begin{qoptions}
\opt{A}{P $\rightarrow$ 4;Q $\rightarrow$ 1;R $\rightarrow$ 5;S $\rightarrow$ 2}
\opt{B}{P $\rightarrow$ 4;Q $\rightarrow$ 1;R $\rightarrow$ 2;S $\rightarrow$ 5}
\opt{C}{P $\rightarrow$ 4;Q $\rightarrow$ 2;R $\rightarrow$ 5;S $\rightarrow$ 3}
\opt{D}{P $\rightarrow$ 1;Q $\rightarrow$ 4;R $\rightarrow$ 2;S $\rightarrow$ 5}
\end{qoptions}
\end{bankq}
\begin{bankpassage}[id=SA23-P1-CHEM-S3-G5-6, source={Score Advanced - Test 23 (14 Apr 2024)}, questions={5-6}]
For the reaction 3A $\rightarrow$ 2B, the following graph is obtained experimentally.\par
\bfig{131.5}{SA23_p1_chem_s3_para5_6_1.png}
\end{bankpassage}
\begin{bankq}[id=SA23-P1-CHEM-S3-Q05, source={Score Advanced - Test 23 (14 Apr 2024)}, paper=1, qno=5, type=numerical, subtopic={Order}, answer={4}, answer_src=key, difficulty=7, diff_parts={K3 S1 I2 C1 F1}, diff_note={concepts: order of reaction from a graph, slope of the log plot | steps: slope of the log plot}, group=SA23-P1-CHEM-S3-G5-6]
What is order of reaction.
\end{bankq}
\begin{bankq}[id=SA23-P1-CHEM-S3-Q06, source={Score Advanced - Test 23 (14 Apr 2024)}, paper=1, qno=6, type=numerical, subtopic={Order}, answer={16}, answer_src=key, difficulty=8, diff_parts={K3 S1 I2 C2 F1}, diff_note={concepts: rate at a given concentration for a 3A to 2B reaction, rate laws, rate constant and order from the graph | steps: rate constant and order from the graph; rate of reaction}, group=SA23-P1-CHEM-S3-G5-6]
The rate of reaction when [A] = 0.2 M is P $\times$ 10$^{-4}$ M min$^{-1}$, the value of `9 P' is.
\end{bankq}
\begin{bankq}[id=SA24-P2-CHEM-S1-Q03, source={Score Advanced - Test 24 (17 Apr 2024)}, paper=2, qno=3, type=single, subtopic={Arrhenius equation}, answer={B}, answer_src=key, difficulty=4, diff_parts={K2 S1 I1 C1 F0}, diff_note={concepts: reason for the doubling of rate on a 10 K rise, arrhenius equation | steps: number of molecules with energy above threshold}]
For most reactions, rate of a reaction increases about two times when its temperature is raised by 10 K. This fact is due to
\begin{qoptions}
\opt{A}{increase in number of collisions per unit volume}
\opt{B}{increase in number of molecules having energy $E_A$ or more}
\opt{C}{lowering of activation energy}
\opt{D}{lowering of pre-exponential factor in Arrhenius equation}
\end{qoptions}
\end{bankq}
\begin{bankpassage}[id=SA25-P2-CHEM-S3-G1-2, source={Score Advanced - Test 25 (21 Apr 2024)}, questions={1-2}]
For the reaction, CHCl$_3$(g) + Cl$_2$(g) $\xrightarrow{h\nu}$ CCl$_4$(g) + HCl(g) mechanism is given below\par
Cl$_2$(g) $\underset{k_{-1}}{\overset{k_1}{\rightleftharpoons}}$ 2Cl(g)\par
CHCl$_3$(g) + Cl(g) $\xrightarrow{k_2}$ CCl$_3$(g) + HCl(g)\par
CCl$_3$(g) + HCl(g) $\xrightarrow{k_3}$ CCl$_4$(g)\par
\textbf{Given :}\par
$k_1=4.8\times10^3$ sec$^{-1}$, $k_{-1}=1.2\times10^3$ M$^{-1}$ sec$^{-1}$,\par
$k_2=1.3\times10^{-2}$ M$^{-1}$ sec$^{-1}$, $k_3=2.1\times10^2$ M$^{-1}$ sec$^{-1}$\par
$E_{a_1}$ = 10 kJ/mol, $E_{a_{-1}}$ = 15 kJ/mol,\par
$E_{a_2}$ = 45 kJ/mol, $E_{a_3}$ = 12 kJ/mol
\end{bankpassage}
\begin{bankq}[id=SA25-P2-CHEM-S3-Q01, source={Score Advanced - Test 25 (21 Apr 2024)}, paper=2, qno=1, type=numerical, subtopic={Rate laws}, answer={2.60}, answer_src=key, difficulty=8, diff_parts={K3 S1 I2 C2 F1}, diff_note={concepts: overall rate constant for a photochemical chain mechanism, rate laws, steady state on Cl radicals | steps: steady state on Cl radicals; combine constants}, group=SA25-P2-CHEM-S3-G1-2]
Overall rate constant of reaction, CHCl$_3$(g) + Cl$_2$(g) $\xrightarrow{h\nu}$ CCl$_4$(g) + HCl(g), is \_\_\_\_\_\_ $\times10^{-2}$ M$^{-1/2}$ sec$^{-1}$.
\end{bankq}
\begin{bankq}[id=SA25-P2-CHEM-S3-Q02, source={Score Advanced - Test 25 (21 Apr 2024)}, paper=2, qno=2, type=numerical, subtopic={Rate laws}, answer={42.50}, answer_src=key, difficulty=8, diff_parts={K3 S1 I2 C2 F1}, diff_note={concepts: overall activation energy of the chain mechanism, arrhenius equation, combination of activation energies of the steps | steps: combination of activation energies of the steps}, group=SA25-P2-CHEM-S3-G1-2]
Overall activation energy of reaction, CHCl$_3$(g) + Cl$_2$(g) $\xrightarrow{h\nu}$ CCl$_4$(g) + HCl(g), is \_\_\_\_\_\_\_\_ kJ/mole.
\end{bankq}
\begin{bankpassage}[id=SA26-P2-CHEM-S3-G1-2, source={Score Advanced - Test 26 (21 Apr 2024)}, questions={1-2}]
For the reaction, CHCl$_3$(g) + Cl$_2$(g) $\xrightarrow{h\nu}$ CCl$_4$(g) + HCl(g) mechanism is given below\par
Cl$_2$(g) $\underset{k_{-1}}{\overset{k_1}{\rightleftharpoons}}$ 2Cl(g)\par
CHCl$_3$(g) + Cl(g) $\xrightarrow{k_2}$ CCl$_3$(g) + HCl(g)\par
CCl$_3$(g) + HCl(g) $\xrightarrow{k_3}$ CCl$_4$(g)\par
\textbf{Given :}\par
$k_1=4.8\times10^3$ sec$^{-1}$, $k_{-1}=1.2\times10^3$ M$^{-1}$ sec$^{-1}$,\par
$k_2=1.3\times10^{-2}$ M$^{-1}$ sec$^{-1}$, $k_3=2.1\times10^2$ M$^{-1}$ sec$^{-1}$\par
$E_{a_1}$ = 10 kJ/mol, $E_{a_{-1}}$ = 15 kJ/mol,\par
$E_{a_2}$ = 45 kJ/mol, $E_{a_3}$ = 12 kJ/mol
\end{bankpassage}
\begin{bankq}[id=SA26-P2-CHEM-S3-Q01, source={Score Advanced - Test 26 (21 Apr 2024)}, paper=2, qno=1, type=numerical, subtopic={Rate laws}, answer={2.60}, answer_src=key, difficulty=8, diff_parts={K3 S1 I2 C2 F1}, diff_note={concepts: overall rate constant for a photochemical chain mechanism, rate laws, steady state on Cl radicals | steps: steady state on Cl radicals; combine constants}, group=SA26-P2-CHEM-S3-G1-2]
Overall rate constant of reaction, CHCl$_3$(g) + Cl$_2$(g) $\xrightarrow{h\nu}$ CCl$_4$(g) + HCl(g), is \_\_\_\_\_\_ $\times10^{-2}$ M$^{-1/2}$ sec$^{-1}$.
\end{bankq}
\begin{bankq}[id=SA26-P2-CHEM-S3-Q02, source={Score Advanced - Test 26 (21 Apr 2024)}, paper=2, qno=2, type=numerical, subtopic={Rate laws}, answer={42.50}, answer_src=key, difficulty=8, diff_parts={K3 S1 I2 C2 F1}, diff_note={concepts: overall activation energy of the chain mechanism, arrhenius equation, combination of activation energies of the steps | steps: combination of activation energies of the steps}, group=SA26-P2-CHEM-S3-G1-2]
Overall activation energy of reaction, CHCl$_3$(g) + Cl$_2$(g) $\xrightarrow{h\nu}$ CCl$_4$(g) + HCl(g), is \_\_\_\_\_\_\_\_ kJ/mole.
\end{bankq}
\begin{bankq}[id=SA28-P1-CHEM-S1-Q04, source={Score Advanced - Test 28 (27 Apr 2024)}, paper=1, qno=4, type=numerical, subtopic={Rate laws}, answer={45.15}, answer_src=key, difficulty=8, diff_parts={K3 S1 I2 C2 F1}, diff_note={concepts: time for a first order decomposition of dinitrogen tetroxide, rate laws, integrated first order law | steps: integrated first order law}]
The rate constant for the reaction N$_2$O$_4$(g) $\longrightarrow$ 2NO$_2$(g) is 92.12 $\times$ 10$^{-2}$ min$^{-1}$. The time (in seconds) required for the molar ratio of NO$_2$ (g) to N$_2$O$_4$ (g) in the reaction mixture to become 2 : 1 is\par
[log 2 = 0.3010, ln 10 = 2.303]
\end{bankq}
\begin{bankq}[id=SA28-P2-CHEM-S1-Q06, source={Score Advanced - Test 28 (27 Apr 2024)}, paper=2, qno=6, type=multi, subtopic={Arrhenius equation}, answer={B,C,D}, answer_src=key, difficulty=7, diff_parts={K3 S1 I2 C1 F1}, diff_note={concepts: statements on the Arrhenius equation, exp(-Ea/RT) is the fraction with energy at least Ea, pre-exponential factor | steps: exp(-Ea/RT) is the fraction with energy at least Ea; pre-exponential factor}]
According to the Arrhenius equation:
\begin{qoptions}
\opt{A}{In Arrhenius equation, factor e$^{-\mathrm{Ea/RT}}$ corresponds to fraction of reactant molecules that have kinetic energy equal to activation energy.}
\opt{B}{Rate constant increases with increases in temperature. This is mainly due to a greater number of collisions whose energy exceeds the activation energy.}
\opt{C}{Higher the magnitude of activation energy of a reaction, stronger is the temperature dependence of its rate constant.}
\opt{D}{The pre$-$exponential factor depends on the frequency at which collisions between reactant molecules per unit volume per unit time occur, irrespective of their energy.}
\end{qoptions}
\end{bankq}
\begin{bankq}[id=SA33-P1-CHEM-S3-Q09, source={Score Advanced - Test 33 (12 May 2024)}, paper=1, qno=9, type=single, subtopic={Rate laws}, answer={B}, answer_src=key, difficulty=4, diff_parts={K2 S1 I1 C1 F0}, diff_note={concepts: standard problem in this chapter, rate laws | steps: apply the relevant results; compute}, also_in={Score Advanced - Test 37 (19 May 2024) P1 Q9}]
Match the half-life of the reaction in List-I with the integrated rate equation in List-II\par
(a = initial concentration of reactant and a $-$ x $-$ concentration of reactant at time `t'.)\par
{\footnotesize\setlength{\tabcolsep}{1.5pt}\noindent\begin{tabular}{|c|L{2.3cm}|c|L{3.6cm}|}
\hline
 & \multicolumn{1}{>{\centering\arraybackslash}p{2.3cm}|}{List $-$ I} & & \multicolumn{1}{>{\centering\arraybackslash}p{3.6cm}|}{List $-$ II} \\ \hline
(P) & $t_{1/2}=\dfrac{0.586}{k}\sqrt a$ & (1) & $k=\dfrac xt$ \\ \hline
(Q) & $t_{1/2}=\dfrac{0.693}{k}$ & (2) & $k=\dfrac1t\left[\dfrac1{a-x}-\dfrac1a\right]$ \\ \hline
(R) & $t_{1/2}=\dfrac a{2k}$ & (3) & $k=\dfrac2t\left[\sqrt a-\sqrt{a-x}\right]$ \\ \hline
(S) & $t_{1/2}=\dfrac1{ka}$ & (4) & $k=\dfrac{2.303}t\log\dfrac a{a-x}$ \\ \hline
 & & (5) & $k=\dfrac2t\left[\dfrac1{\sqrt{a-x}}-\dfrac1{\sqrt a}\right]$ \\ \hline
\end{tabular}}\par
\begin{qoptions}
\opt{A}{P $\to$ 5;Q $\to$ 4;R $\to$ 1;S $\to$ 2}
\opt{B}{P $\to$ 3;Q $\to$ 4;R $\to$ 1;S $\to$ 2}
\opt{C}{P $\to$ 2;Q $\to$ 3;R $\to$ 4;S $\to$ 2}
\opt{D}{P $\to$ 3;Q $\to$ 2;R $\to$ 4;S $\to$ 5}
\end{qoptions}
\end{bankq}
\begin{bankq}[id=SA33-P2-CHEM-S1-Q02, source={Score Advanced - Test 33 (12 May 2024)}, paper=2, qno=2, type=single, subtopic={Rate laws}, answer={A}, answer_src=key, difficulty=4, diff_parts={K2 S1 I1 C1 F0}, diff_note={concepts: standard problem in this chapter, rate laws | steps: apply the relevant results; compute}, also_in={Score Advanced - Test 34 (12 May 2024) P2 Q2}]
Inversion of sucrose (C$_{12}$H$_{22}$O$_{11}$) is first-order reaction and is studied by measuring angle of rotation of the reaction mixture at different instants of time\par
C$_{12}$H$_{22}$O$_{11}$ + H$_2$O $\xrightarrow{\mathrm H^+}$ C$_6$H$_{12}$O$_6$ + C$_6$H$_{12}$O$_6$\par
If r$_0$, r$_\mathrm t$ and r$_\infty$ are the angles of rotation of the reaction mixture at the start, at the time t and at the end of the reaction respectively, then 50\% inversion of sucrose occurs when :
\begin{qoptions}
\opt{A}{r$_0$ = 2r$_\mathrm t$ $-$ r$_\infty$}
\opt{B}{r$_0$ = r$_\mathrm t$ $-$ r$_\infty$}
\opt{C}{r$_0$ = r$_\mathrm t$ $-$ 2r$_\infty$}
\opt{D}{r$_0$ = r$_\mathrm t$ + r$_\infty$}
\end{qoptions}
\end{bankq}
\begin{bankq}[id=SA35-P1-CHEM-S2-Q09, source={Score Advanced - Test 35 (15 May 2024)}, paper=1, qno=9, type=single, subtopic={Rate laws}, answer={A}, answer_src=key, difficulty=8, diff_parts={K3 S2 I2 C2 F0}, diff_note={concepts: matching kinetics quantities, rate laws, stoichiometric ratios | steps: stoichiometric ratios; half life order; first order time; Nernst}]
Match the list\par
{\footnotesize\setlength{\tabcolsep}{1.5pt}\noindent\begin{tabular}{|c|L{4.6cm}|c|c|}
\hline
 & \multicolumn{1}{c|}{List-I} & \multicolumn{2}{c|}{List-II} \\ \hline
(P) & Ratio of rate of appearance (gm/min) of SO$_3$(g) and rate of disappearance (gm/min) of SO$_2$(g) for : 2SO$_2$(g) + O$_2$(g) $\longrightarrow$ 2SO$_3$(g) ; would be & (1) & 0 \\ \hline
(Q) & Based on following the observations for a reaction A $\xrightarrow{k}$ product\newline [A]$_0$ (in M) : 0.0677 \ \ 0.136 \ \ 0.272\newline t$_{1/2}$ (in min) : 2 \ \ 4 \ \ 8\newline the order of reaction would be & (2) & 3 \\ \hline
(R) & Time (min) taken by a reaction to decrease the concentration of the reactant(R) from 0.2 mole/litre to 0.025 mole/litre if its rate constant corresponding to rate of disappearance of R is 0.693 min$^{-1}$, would be & (3) & 1.25 \\ \hline
(S) & If emf at 25$^\circ$C of the following Galvanic cell\newline Pt $|$ H$_2$(g) $|$ CH$_3$COOH(aq.) $\|$ HCOOH(aq.) $|$ H$_2$(g) $|$ Pt\newline \hspace*{4.6em}(C$_1$)\hspace*{4.2em}(C$_2$)\newline is 0.0177 volt then ratio of [H$^+$] in anodic compartment and cathodic compartment would be approximately & (4) & 2 \\ \hline
\end{tabular}}\par
\begin{qoptions}
\opt{A}{P $\to$ 3;Q $\to$ 1;R $\to$ 2;S $\to$ 4}
\opt{B}{P $\to$ 2;Q $\to$ 2;R $\to$ 3;S $\to$ 1}
\opt{C}{P $\to$ 1;Q $\to$ 3;R $\to$ 3;S $\to$ 4}
\opt{D}{P $\to$ 3;Q $\to$ 2;R $\to$ 1;S $\to$ 4}
\end{qoptions}
\end{bankq}
\begin{bankq}[id=SA38-P2-CHEM-S4-Q08, source={Score Advanced - Test 38 (13 Apr 2025)}, paper=2, qno=8, type=numerical, subtopic={Nuclear chemistry}, answer={180}, answer_src=key, difficulty=8, diff_parts={K3 S1 I2 C2 F1}, diff_note={concepts: standard problem in this chapter, nuclear chemistry, apply the relevant results | steps: apply the relevant results; compute}]
A sample initially contains only U-238 isotope of uranium. With time, some of the U-238 radioactively decays into Pb-206 while the rest of it remains undisintegrated.\par
When the age of the sample is \textbf P $\times$ 10$^8$ years, the ratio of moles of Pb-206 to that of U-238 in the sample is found to be 15. The value of \textbf P is \underline{\hspace{1.5cm}}.\par
[Given : Half-life of U-238 is 4.5 $\times$ 10$^9$ years]
\end{bankq}
\begin{bankq}[id=SA40-P1-CHEM-S4-Q05, source={Score Advanced - Test 40 (27 Apr 2025)}, paper=1, qno=5, type=numerical, subtopic={Order}, answer={2}, answer_src=key, difficulty=8, diff_parts={K3 S1 I2 C2 F1}, diff_note={concepts: standard problem in this chapter, order, apply the relevant results | steps: apply the relevant results; compute}]
If reaction A $\xrightarrow{\mathrm K}$ B has K = 6.93 $\times$ 10$^{-3}$ min$^{-1}$ at 20$^\circ$C, then half-life of `A' at 70$^\circ$C is $\dfrac{25\mathrm x}{16}$ min. Find value of `x', if temperature coefficient of reaction is assumed to have a constant value of 2 ?
\end{bankq}
\begin{bankq}[id=SA40-P1-CHEM-S4-Q06, source={Score Advanced - Test 40 (27 Apr 2025)}, paper=1, qno=6, type=numerical, subtopic={Order}, answer={5}, answer_src=key, difficulty=8, diff_parts={K3 S1 I2 C2 F1}, diff_note={concepts: standard problem in this chapter, order, apply the relevant results | steps: apply the relevant results; compute}]
SO$_3$ decomposes into SO$_2$ \& O$_2$ according to this equation in a rigid closed vessel of volume 1 litre.\par
SO$_3$(g)$\longrightarrow$SO$_2$(g) + $\dfrac12$O$_2$(g)\par
At time = 20 min, rate of appearance of O$_2$ was 20 g L$^{-1}$ per minute and $\dfrac{\mathrm d[\mathrm{SO}_3]}{\mathrm{dt}}$ at that instant is reported as $\dfrac{-\mathrm x}3$g L$^{-1}$s$^{-1}$. Then value of `x' is? [Atomic mass of S = 32]
\end{bankq}
\begin{bankpassage}[id=SA40-P2-CHEM-S2-G3-4, source={Score Advanced - Test 40 (27 Apr 2025)}, questions={3-4}]
The vapour pressure of two miscible liquids A and B are 300 and 500 mm Hg respectively. In a flask, 10 mole of A is mixed with 12 mole of B. However, as soon as `B' is added, `A' starts polymerizing into a completely insoluble solid. The polymerisation follows first order kinetics. After 100 min, 1 mole of a non$-$volatile solute is dissolved, which arrests the polymerisation completely. The final vapour pressure of the solution is 400 mm of Hg.\par
Assume negligible volume change on mixing and polymerisation and ideal behaviour for the final solution. [log2 = 0.3, log$_\mathrm e$10 = 2.3]
\end{bankpassage}
\begin{bankq}[id=SA40-P2-CHEM-S2-Q03, source={Score Advanced - Test 40 (27 Apr 2025)}, paper=2, qno=3, type=numerical, subtopic={Rate laws}, answer={10}, answer_src=key, difficulty=8, diff_parts={K3 S1 I2 C2 F1}, diff_note={concepts: standard problem in this chapter, rate laws, apply the relevant results | steps: apply the relevant results; compute}, group=SA40-P2-CHEM-S2-G3-4]
If rate constant for the first order is `K' min$^{-1}$, then find the value of $\left(\dfrac{10^4\mathrm K}{2.3}\right)$ is:
\end{bankq}
\begin{bankq}[id=SA40-P2-CHEM-S2-Q04, source={Score Advanced - Test 40 (27 Apr 2025)}, paper=2, qno=4, type=numerical, subtopic={Rate laws}, answer={4}, answer_src=solved, working={Negative deviation: i, ii, iii (mixing), v true; dS\_surr>0 since exothermic so iv false: 4.}, difficulty=8, diff_parts={K3 S1 I2 C2 F1}, diff_note={concepts: standard problem in this chapter, raoult's law, apply the relevant results | steps: apply the relevant results; compute}, group=SA40-P2-CHEM-S2-G3-4]
If initially mixed A and B has non$-$ideal behaviour with negative deviation from Raoult's law and there is no polymerisation, then how many of the following are true:\par
(i) Attraction between A and B is more than individual attraction of A.....A and B.....B.\par
(ii) Potential energy of the system decreases on mixing.\par
(iii) $\Delta$S$_{\mathrm{System}}$ $>$ 0\par
(iv) $\Delta$S$_{\mathrm{Surrounding}}$ $<$ 0\par
(v) $\Delta$G$_{\mathrm{Mixing}}$ $<$ 0
\end{bankq}
\begin{bankq}[id=SA41-P1-CHEM-S3-Q10, source={Score Advanced - Test 41}, paper=1, qno=10, type=single, subtopic={Arrhenius equation}, answer={B}, answer_src=key, difficulty=4, diff_parts={K2 S1 I1 C1 F0}, diff_note={concepts: standard problem in this chapter, arrhenius equation | steps: apply the relevant results; compute}]
t$_\mathrm{x/y}$ = time in which x/y fraction of reactant converts into product :\par
{\footnotesize\setlength{\tabcolsep}{1.5pt}\noindent\begin{tabular}{|c|L{0.9cm}|c|L{4.9cm}|}\hline & \multicolumn{1}{>{\centering\arraybackslash}p{0.9cm}|}{List-I} & & \multicolumn{1}{>{\centering\arraybackslash}p{4.9cm}|}{List-II} \\ \hline
(P) & t$_{5/9}$ & (1) & Equal to 54 sec. if t$_{1/3}$ is 18 sec. in case of first order reaction. \\ \hline
(Q) & t$_{19/27}$ & (2) & Equal to 32 sec. if t$_{1/4}$ is 16 sec. in case of first order reaction. \\ \hline
(R) & t$_{7/8}$ & (3) & Equal to 56 sec. if t$_{1/3}$ is 4 sec. in case of second order reaction. \\ \hline
(S) & t$_{7/16}$ & (4) & Equal to 30 sec. if t$_{1/3}$ is 18 sec. in case of zero order reaction. \\ \hline
 & & (5) & Equal to 28 sec. if t$_{1/2}$ is 16 sec. in case of zero order reaction. \\ \hline
\end{tabular}}\par
\begin{qoptions}
\opt{A}{P $\to$ 1;Q $\to$ 2;R $\to$ 3;S $\to$ 4,5}
\opt{B}{P $\to$ 4;Q $\to$ 1;R $\to$ 3,5;S $\to$ 2}
\opt{C}{P $\to$ 2;Q $\to$ 3;R $\to$ 1,3;S $\to$ 5}
\opt{D}{P $\to$ 1,3;Q $\to$ 2,4;R $\to$ 3;S $\to$ 5}
\end{qoptions}
\end{bankq}
\begin{bankq}[id=SA41-P1-CHEM-S4-Q03, source={Score Advanced - Test 41}, paper=1, qno=3, type=single, subtopic={Arrhenius equation}, answer={8500}, answer_src=key, difficulty=7, diff_parts={K3 S1 I2 C2 F0}, diff_note={concepts: standard problem in this chapter, arrhenius equation, apply the relevant results | steps: apply the relevant results; compute}]
Consider the following reversible reaction,\par
A(g) + B(g) $\rightleftharpoons$ AB(g)\par
The activation energy of the backward reaction exceeds that of the forward reaction by 2RT (in J mol$^{-1}$). If the pre-exponential factor of the forward reaction is 4 times that of the reverse reaction, the absolute value of $\Delta$G$^\theta$ (in J mol$^{-1}$) for the reaction at 300 K is .........\par
(Given : ln (2) = 0.7, RT = 2500 J mol$^{-1}$ at 300 K and G is the Gibbs energy)
\end{bankq}

\section*{SRGBT papers}
\begin{bankq}[id=SRGBT12-P1-CHEM-S1-Q30, source={SRGBT - Test 12 (10 Oct 2021)}, paper=1, qno=30, type=multi, subtopic={Rate laws}, answer={A,B}, answer_src=key, difficulty=8, diff_parts={K3 S1 I2 C2 F1}, diff_note={concepts: standard problem in this chapter, rate laws, apply the relevant results | steps: apply the relevant results; compute}]
For the reaction N$_2$O$_5$(g) $\longrightarrow$ 2NO$_2$(g) + $\dfrac12$O$_2$(g)\par
N$_2$O$_5$(g) decomposed at a constant volume and temperature. Initial pressure of N$_2$O$_5$ is 600 mm of Hg and total pressure at equilibrium is 960 mm of Hg. Assuming ideal behavior of gases, Select the correct option(s).
\begin{qoptions}
\opt{A}{Involved reaction is intramolecular redox reaction}
\opt{B}{The mole fraction of NO$_2$ in the equilibrium mixture is 0.5}
\opt{C}{The equilibrium constant K$_\mathrm p$ for the above reaction is 540 mm of Hg}
\opt{D}{For the above reaction K$_\mathrm C$ = K$_\mathrm p$ (RT)$^3$}
\end{qoptions}
\end{bankq}

\section*{Enthusiast Pre-Score papers}
\begin{bankq}[id=EPS03-P1-CHEM-S1-Q03, source={Enthusiast Pre-Score - Test 03 (23 Nov 2020)}, paper=1, qno=3, type=multi, subtopic={Order}, answer={A,B,C,D}, answer_src=key, difficulty=7, diff_parts={K2 S2 I2 C1 F1}, diff_note={concepts: zero order reaction with stoichiometry, rate constant from data | steps: rate constant from data; average time; times for fractions}]
For a zero order reaction: A $\to$ 2B, [A] = 0.4M and [B] = 0 at t = 0. If at t = 25 min, [B] = 0.5M, then the correct information(s) about the reaction, with respect to initial condition, is/are:
\begin{qoptions}
\opt{A}{t$_{100\%}$ = 40 min}
\opt{B}{T$_{av}$ = 20 min}
\opt{C}{t$_{50\%}$ = 20 min}
\opt{D}{t$_{10\%}$ = 4 min}
\end{qoptions}
\end{bankq}
\begin{bankq}[id=EPS03-P1-CHEM-S1-Q04, source={Enthusiast Pre-Score - Test 03 (23 Nov 2020)}, paper=1, qno=4, type=multi, subtopic={Nuclear chemistry}, answer={B,D}, answer_src=key, difficulty=4, diff_parts={K1 S2 I1 C0 F1}, diff_note={concepts: nuclear processes | steps: decay modes; conversions of nucleons; check statements}]
The correct information(s) about nuclear processes is/are:
\begin{qoptions}
\opt{A}{In alpha decay, two protons and two neutrons come out from the nucleus, in succession.}
\opt{B}{An independent neutron may convert spontaneously into proton.}
\opt{C}{An independent proton may convert spontaneously into neutron.}
\opt{D}{In a positron decay, a proton spontaneously converts into neutron.}
\end{qoptions}
\end{bankq}
\begin{bankq}[id=EPS03-P2-CHEM-S2-Q07, source={Enthusiast Pre-Score - Test 03 (23 Nov 2020)}, paper=2, qno=7, type=single, subtopic={Nuclear chemistry}, answer={A}, answer_src=key, difficulty=3, diff_parts={K1 S2 I1 C0 F0}, diff_note={concepts: decay modes and neutron proton ratio | steps: ratio relative to the belt; mode of decay; match}]
Match correctly the nature of radioactive isotopes in List-II with their decay mode in List-I.\par
\noindent{\footnotesize\setlength{\tabcolsep}{2pt}\begin{tabular}{|L{.8cm}|L{3cm}|L{.8cm}|L{2.8cm}|}\hline
\multicolumn{2}{|c|}{List-I} & \multicolumn{2}{c|}{List-II}\\\hline
(P) & Alpha decay & (1) & Al$^{24}$\\\hline
(Q) & Electron (beta decay) & (2) & Al$^{29}$\\\hline
(R) & Positron decay & (3) & Na$^{20}$\\\hline
(S) & Gamma decay & (4) & U$^{238}$\\\hline
 & & (5) & Mg$^{26}$\\\hline
\end{tabular}}
\begin{qoptions}
\opt{A}{P $\to$ 4; Q $\to$ 2, 5; R $\to$ 1, 3 ; S $\to$ 1, 2, 3, 4, 5}
\opt{B}{P $\to$ 4; Q $\to$ 1, 3; R $\to$ 2, 5; S $\to$ 1, 2, 3, 4, 5}
\opt{C}{P $\to$ 4; Q $\to$ 2, 3; R $\to$ 1, 5; S $\to$ 1, 2, 3, 4, 5}
\opt{D}{P $\to$ 4; Q $\to$ 1, 5; R $\to$ 2, 5 ; S $\to$ 1, 2, 3, 4, 5}
\end{qoptions}
\end{bankq}
\begin{bankq}[id=EPS03-P2-CHEM-S3-Q04, source={Enthusiast Pre-Score - Test 03 (23 Nov 2020)}, paper=2, qno=4, type=numerical, subtopic={Order}, answer={10.00}, answer_src=key, difficulty=8, diff_parts={K2 S2 I2 C2 F1}, diff_note={concepts: first order gaseous reaction by pressure rise, pressure relation | steps: pressure relation; rate constant; half life}]
For a first order gaseous reaction occurring at constant volume and temperature aA $\to$ (b - 1) B + C (with b $>$ a), the pressure of the system rose by $50\left(\dfrac ba-1\right)$ \% in 10 min. The half life of reaction (in min) is.
\end{bankq}

\section*{TOAS papers}
\begin{bankq}[id=TOAS04-P1-CHEM-S1-Q02, source={TOAS - Test 04 (15 Mar 2022)}, paper=1, qno=2, type=single, subtopic={Rate laws}, answer={B}, answer_src=key, difficulty=1, diff_parts={K1 S1 I0 C0 F0}, diff_note={concepts: rate law from initial rates | steps: order in each reactant from the table; rate constant and the requested rate}]
Initial rate data at $25^\circ$C are listed in the table for the reaction

$\mathrm{NH_{4(aq)}^{+}+NO_{2\,(aq)}^{-}\to N_{2(g)}+2H_2O_{(l)}}$
{\small\setlength{\tabcolsep}{3pt}\renewcommand{\arraystretch}{1.25}\begin{center}\begin{tabular}{|c|c|c|c|}
\hline
Expt.\ No. & Initial $\left[\mathrm{NH_4^+}\right]$ & Initial $\left[\mathrm{NO_2^-}\right]$ & \begin{tabular}[c]{@{}c@{}}Initial rate of\\ consumption of\\ $\mathrm{NH_4^+}$ M/s\end{tabular}\\
\hline
1. & 0.24 & 0.10 & $7.2\times10^{-6}$\\
\hline
2. & 0.12 & 0.10 & $3.6\times10^{-6}$\\
\hline
3. & 0.12 & 0.15 & $5.4\times10^{-6}$\\
\hline
\end{tabular}\end{center}}
What is the reaction rate when concentrations are $\left[\mathrm{NH_4^+}\right]=0.39$ M and $\left[\mathrm{NO_2^-}\right]=0.052$ M ?
\begin{qoptions}
\opt{A}{$6.1\times10^{-5}$ M/s}
\opt{B}{$6.1\times10^{-6}$ M/s}
\opt{C}{$6.1\times10^{-7}$ M/s}
\opt{D}{$3\times10^{-5}$ M/s}
\end{qoptions}
\end{bankq}
\begin{bankpassage}[id=TOAS04-P1-CHEM-S3-G3-4, source={TOAS - Test 04 (15 Mar 2022)}, questions={3-4}]
For the first order dissociation reaction

$\mathrm{A_2B_3(aq.)\to2A^{3+}(aq.)+3B^{2-}(aq.)}$

a definite volume of solution is taken at different instant and its osmotic pressure rise is measured. Osmotic pressure rise at $t=0$ is 2 mm and at $t=10$ min is 6 mm. [Given density of solution is constant and equal to 1 gm/ml, $g=10\ \mathrm{m/s^2}$]
\end{bankpassage}
\begin{bankq}[id=TOAS04-P1-CHEM-S3-Q03, source={TOAS - Test 04 (15 Mar 2022)}, paper=1, qno=3, type=numerical, subtopic={Order}, answer={8.00}, answer_src=key, difficulty=3, diff_parts={K2 S1 I0 C0 F1}, diff_note={concepts: first order decay, osmotic pressure | steps: concentration at 20 min; osmotic pressure}, group=TOAS04-P1-CHEM-S3-G3-4]
Calculate osmotic pressure of solution in mm at $t=20$ min.
\end{bankq}
\begin{bankq}[id=TOAS04-P1-CHEM-S3-Q04, source={TOAS - Test 04 (15 Mar 2022)}, paper=1, qno=4, type=numerical, subtopic={Order}, answer={30.00}, answer_src=key, difficulty=2, diff_parts={K1 S1 I0 C0 F1}, diff_note={concepts: first order kinetics | steps: 87.5 percent completion is three half-lives}, group=TOAS04-P1-CHEM-S3-G3-4]
Calculate time (in minute) for 87.5 \% of the reaction to complete -
\end{bankq}
\begin{bankq}[id=TOAS04-P1-CHEM-S4-Q08, source={TOAS - Test 04 (15 Mar 2022)}, paper=1, qno=8, type=numerical, subtopic={Arrhenius equation}, answer={4}, answer_src=key, difficulty=3, diff_parts={K2 S1 I0 C0 F1}, diff_note={concepts: parallel first order kinetics, Arrhenius equation | steps: rate constant ratios from the product ratio; net activation energy as the weighted average}]
A substance undergoes first order parallel reaction as shown. Calculate net activation energy (in kcal) if it is a known that $B$, $C$ and $D$ are formed in a molar ratio of $2:3:4$ respectively.
\bfig{224.4}{TOAS04_p1_chem_s4_q08_1.png}
\end{bankq}
\begin{bankq}[id=TOAS04-P1-CHEM-S4-Q09, source={TOAS - Test 04 (15 Mar 2022)}, paper=1, qno=9, type=numerical, subtopic={Order}, answer={64}, answer_src=key, difficulty=3, diff_parts={K1 S1 I1 C0 F1}, diff_note={concepts: order from half lives | steps: half life scaling t proportional to a\^{}(1-n) gives n = 4; half life at the new concentration | I: scaling of half life}, also_in={TOAS - Test 09 (30 May 2022) P2 Q8}]
For a reaction $A\to P$

Half lite measured for two different values of initial concentrations $5\times10^{-3}$ M and $25\times10^{-4}$ M are 1 and 8 hrs respectively. If initial concentration is adjusted to $1.25\times10^{-3}$ M, what is the half-life? (in hrs)
\end{bankq}
\begin{bankq}[id=TOAS04-P2-CHEM-S1-Q03, source={TOAS - Test 04 (15 Mar 2022)}, paper=2, qno=3, type=multi, subtopic={Rate laws}, answer={A,C,D}, answer_src=key, difficulty=4, diff_parts={K1 S2 I0 C1 F1}, diff_note={concepts: first order gaseous decomposition | steps: total pressure at infinite time; at 10 percent decomposition; at 5 percent decomposition; and the half life}, also_in={TOAS - Test 12 (10 Jul 2022) P1 Q3}]
In an evacuated vessel of 8.2 L capacity, 0.2 moles of $\mathrm{A(g)}$ is taken, vessel is sealed and heated to $127^\circ$C, by which the following first order reaction occurred at constant volume and temperature :-\\
$\mathrm{A(g)}\rightarrow\mathrm{2B(g)+C(s)}$\\
The vapour pressure of $\mathrm{C(s)}$ is 0.1 atm at $127^\circ$C. If the half life of reaction is 20 min, then the correct information(s) related with reaction is/are : ($R=0.082$ L-atm/K-mol)
\begin{qoptions}
\opt{A}{$P_{\mathrm{total}}=1.7$ atm\quad at $t=\infty$}
\opt{B}{$P_{\mathrm{total}}=0.98$ atm after 10\% decomposition of $\mathrm{A(g)}$}
\opt{C}{$P_{\mathrm{total}}=0.88$ atm after 5\% decompositon of $\mathrm{A(g)}$}
\opt{D}{$P_{\mathrm{total}}=1.3$ atm at $t=20$ min}
\end{qoptions}
\end{bankq}
\begin{bankq}[id=TOAS05-P1-CHEM-S2-Q09, source={TOAS - Test 05 (03 Apr 2022)}, paper=1, qno=9, type=single, subtopic={Order}, answer={D}, answer_src=key, difficulty=2, diff_parts={K1 S2 I0 C0 F0}, diff_note={concepts: order of reaction and half-life relations | steps: derive t1/2 and t(alpha) dependence for each order; match List-I; check combinations}]
\textbf{Answer the following by approximately matching the lists based on the information given in the paragraph.}\par
For an $n^{\mathrm{th}}$ order reaction $\mathrm A\to$ products,\par
The rate law is $\dfrac{-dC_A}{dt}=kC_A^{\,n}$. Where $C_A$ = concentration of A at time `$t$'. Let $C_{A_0}$ is initial concentration of A and and $\alpha$ = degree of dissociation of reactant `A'
{\small\setlength{\tabcolsep}{3pt}\renewcommand{\arraystretch}{1.25}\begin{center}\begin{tabular}{|L{0.85cm}|L{1.2cm}|L{0.65cm}|L{3.5cm}|}
\hline
 & List-I & & List-II\\
\hline
(I) & $n=0$ & (P) & $\alpha=\dfrac{C_{A_0}-C_A}{C_{A_0}}$\\
\hline
(II) & $n=\dfrac12$ & (Q) & $\alpha=1-e^{-kt}$\\
\hline
(III) & $n=1$ & (R) & $\alpha=\dfrac{kt}{C_{A_0}}$\\
\hline
(IV) & $n=2$ & (S) & $\alpha=\dfrac{C_{A_0}kt}{(1+C_{A_0}kt)}$\\
\hline
 & & (T) & $\alpha$ = independent of initial concentration of reactant\\
\hline
 & & (U) & Graph between `$\alpha$' and `$t$' is a straight line passing through origin\\
\hline
\end{tabular}\end{center}}
Which of the following options has the correct combination considering List-I and List-II ?
\begin{qoptions}
\opt{A}{II $\to$ P, R}
\opt{B}{I $\to$ P,R,T}
\opt{C}{III $\to$ P,Q,U}
\opt{D}{III $\to$ P,Q,T}
\end{qoptions}
\end{bankq}
\begin{bankq}[id=TOAS05-P1-CHEM-S2-Q10, source={TOAS - Test 05 (03 Apr 2022)}, paper=1, qno=10, type=single, subtopic={Order}, answer={C}, answer_src=key, difficulty=2, diff_parts={K1 S2 I0 C0 F0}, diff_note={concepts: order of reaction and half-life relations | steps: derive concentration-time relations for each order; match List-I; check combinations}]
\textbf{Answer the following by approximately matching the lists based on the information given in the paragraph.}\par
For an $n^{\mathrm{th}}$ order reaction $\mathrm A\to$ products,\par
The rate law is $\dfrac{-dC_A}{dt}=kC_A^{\,n}$. Where $C_A$ = concentration of A at time `$t$'. Let $C_{A_0}$ is initial concentration of A and and $\alpha$ = degree of dissociation of reactant `A'
{\small\setlength{\tabcolsep}{3pt}\renewcommand{\arraystretch}{1.25}\begin{center}\begin{tabular}{|L{0.85cm}|L{1.2cm}|L{0.65cm}|L{3.5cm}|}
\hline
 & List-I & & List-II\\
\hline
(I) & $n=0$ & (P) & $\alpha=\dfrac{C_{A_0}-C_A}{C_{A_0}}$\\
\hline
(II) & $n=\dfrac12$ & (Q) & $\alpha=1-e^{-kt}$\\
\hline
(III) & $n=1$ & (R) & $\alpha=\dfrac{kt}{C_{A_0}}$\\
\hline
(IV) & $n=2$ & (S) & $\alpha=\dfrac{C_{A_0}kt}{(1+C_{A_0}kt)}$\\
\hline
 & & (T) & $\alpha$ = independent of initial concentration of reactant\\
\hline
 & & (U) & Graph between `$\alpha$' and `$t$' is a straight line passing through origin\\
\hline
\end{tabular}\end{center}}
Which of the following options has the correct combination considering List-I and List-II ?
\begin{qoptions}
\opt{A}{I $\to$ P,Q}
\opt{B}{IV $\to$ S,T}
\opt{C}{IV $\to$ P,S}
\opt{D}{III $\to$ Q,U}
\end{qoptions}
\end{bankq}
\begin{bankq}[id=TOAS05-P2-CHEM-S1-Q10, source={TOAS - Test 05 (03 Apr 2022)}, paper=2, qno=10, type=single, subtopic={Order}, answer={B}, answer_src=key, difficulty=2, diff_parts={K1 S1 I1 C0 F0}, diff_note={concepts: acid catalysed ester hydrolysis | steps: rate constant proportional to [H+] so sulphuric acid is faster | I: dibasic acid supplies more H+}]
The hydrolysis of an ester was carried out separately with 0.05 M HCl and 0.05 M $\mathrm{H_2SO_4}$. Which of the following will be true ?\\
\relax[$k$ = first order rate constant of reaction]
\begin{qoptions}
\opt{A}{$k_{\mathrm{HCl}}>k_{\mathrm{H_2SO_4}}$}
\opt{B}{$k_{\mathrm{HCl}}<k_{\mathrm{H_2SO_4}}$}
\opt{C}{$k_{\mathrm{HCl}}=k_{\mathrm{H_2SO_4}}$}
\opt{D}{$k_{\mathrm{HCl}}=2k_{\mathrm{H_2SO_4}}$}
\end{qoptions}
\end{bankq}
\begin{bankq}[id=TOAS06-P1-CHEM-S1-Q01, source={TOAS - Test 06 (17 Apr 2022)}, paper=1, qno=1, type=single, subtopic={Order}, answer={B}, answer_src=key, difficulty=3, diff_parts={K1 S2 I1 C0 F0}, diff_note={concepts: first order kinetics from optical rotation | steps: rotation is proportional to [A]; rate constant from the two readings; titration volume proportional to B formed | I: relate B produced to A consumed}]
In a hypothetical reaction, $\mathrm{A(aq)\rightleftharpoons 2B(aq)+C(aq)}$ ($1^{\mathrm{st}}$ order decomposition). `A' is optically active (dextro-rototory) while `B' and `C' are optically inactive but `B' takes part in a titration reaction (fast reaction) with $\mathrm{H_2O_2}$. Hence, the progress of reaction can be monitored by measuring rotation of plane polarised light or by measuring volume of $\mathrm{H_2O_2}$ consumed in titration. In an experiment the optical rotation was found to be $\theta=40^\circ$ at $t=20$ min and $\theta=10^\circ$ at $t=50$ min. from start of the reaction. If the progress would have been monitored by titration method, volume of $\mathrm{H_2O_2}$ consumed at $t=15$ min. (from start) is 40 ml then volume of $\mathrm{H_2O_2}$ consumed at $t=60$ min will be :
\begin{qoptions}
\opt{A}{60 ml}
\opt{B}{75 ml}
\opt{C}{52.5 ml}
\opt{D}{90 ml}
\end{qoptions}
\end{bankq}
\begin{bankq}[id=TOAS07-P1-CHEM-S2-Q08, source={TOAS - Test 07 (01 May 2022)}, paper=1, qno=8, type=multi, subtopic={Rate laws}, answer={A,B}, answer_src=key, difficulty=3, diff_parts={K1 S2 I0 C0 F1}, diff_note={concepts: rate expressions and temperature dependence | steps: rate of disappearance relations; activation energy effects; half-life dependence on concentration}]
Which of the following statements are not correct ?
\begin{qoptions}
\opt{A}{For a reaction $(1/2)\mathrm{A}\to2\mathrm{B}$, the rate of disappearance of A is $-2\,d[\mathrm{A}]/dt$}
\opt{B}{For a reaction $(1/2)\mathrm{A_2}\to2\mathrm{B}$, the rate of disappearance of $\mathrm{A_2}$ is related to the rate of appearance of B through an expression $-d[\mathrm{A_2}]/dt=4\,d[\mathrm{B}]/dt$.}
\opt{C}{Reactions with small activation energies slow a lesser increase in rate constant with rise in temperature as compared to that having high activation energy.}
\opt{D}{If the half-life for a given reaction having single reactant is doubled as the initial concentration of a reactant was doubled then order of the reaction would be zero.}
\end{qoptions}
\end{bankq}
\begin{bankpassage}[id=TOAS07-P1-CHEM-S3-G3-4, source={TOAS - Test 07 (01 May 2022)}, questions={3-4}]
For the zero-order reaction :\\
$\mathrm{4A(g)\to3B(g)+2C(g)}$\\
At $t=0$, partial pressure of A $=4000$ mm Hg\\
Rate constant of reaction $=5$ mm Hg /sec
\end{bankpassage}
\begin{bankq}[id=TOAS07-P1-CHEM-S3-Q03, source={TOAS - Test 07 (01 May 2022)}, paper=1, qno=3, type=numerical, subtopic={Order}, answer={5000.00}, answer_src=key, difficulty=3, diff_parts={K1 S2 I0 C0 F1}, diff_note={concepts: first order gas decomposition | steps: rate constant from the half life; pressure of reactant at 250 s; total pressure with stoichiometry}, group=TOAS07-P1-CHEM-S3-G3-4]
Find the total final pressure of gas mixture at $t=250$ second (in mm Hg)
\end{bankq}
\begin{bankq}[id=TOAS07-P1-CHEM-S3-Q04, source={TOAS - Test 07 (01 May 2022)}, paper=1, qno=4, type=numerical, subtopic={Order}, answer={1500.00}, answer_src=key, difficulty=3, diff_parts={K1 S2 I0 C0 F1}, diff_note={concepts: first order gas decomposition | steps: rate constant; pressure of A at 100 s; pressure of B from stoichiometry}, group=TOAS07-P1-CHEM-S3-G3-4]
Find the partial pressure of `B' at $t=100$ second (in mm Hg)
\end{bankq}
\begin{bankq}[id=TOAS07-P1-CHEM-S4-Q08, source={TOAS - Test 07 (01 May 2022)}, paper=1, qno=8, type=numerical, subtopic={Order}, answer={3}, answer_src=key, difficulty=3, diff_parts={K1 S2 I0 C0 F1}, diff_note={concepts: order of reaction from rate data | steps: order with respect to NO from the ratio; order with respect to H2; overall order}]
For the reaction $\mathrm{2NO+H_2\to N_2O+H_2O}$ the value of $-dp/dt$ was found to be 1.50 Torr s$^{-1}$ for a pressure of 304 Torr of NO and 0.375 Torr s$^{-1}$ for a pressure of 152 Torr of NO, the pressure of $\mathrm{H_2}$ being constant. On the other hand, when the pressure of NO was kept constant, $-dp/dt$ was 1.60 Torr s$^{-1}$ for a hydrogen pressure of 294 Torr and 0.8 Torr s$^{-1}$ for a hydrogen pressure of 147 Torr. Determine the overall order of the reaction.
\end{bankq}
\begin{bankpassage}[id=TOAS07-P2-CHEM-S3-G3-4, source={TOAS - Test 07 (01 May 2022)}, questions={3-4}]
Following plot is drawn for a first order reaction $\mathrm{A\to2B}$ [Given $\ln2=0.7$ ]
\bfig{123.5}{TOAS07_p2_chem_s3_para3_4_1.png}
\end{bankpassage}
\begin{bankq}[id=TOAS07-P2-CHEM-S3-Q03, source={TOAS - Test 07 (01 May 2022)}, paper=2, qno=3, type=numerical, subtopic={Order}, answer={0.56}, answer_src=key, difficulty=2, diff_parts={K1 S1 I0 C0 F1}, diff_note={concepts: half-life from kinetic data | steps: rate constant from the data; half-life}, group=TOAS07-P2-CHEM-S3-G3-4]
What is the half life of reaction in minutes ?
\end{bankq}
\begin{bankq}[id=TOAS07-P2-CHEM-S3-Q04, source={TOAS - Test 07 (01 May 2022)}, paper=2, qno=4, type=numerical, subtopic={Order}, answer={1.25}, answer_src=key, difficulty=2, diff_parts={K1 S1 I0 C0 F1}, diff_note={concepts: initial rate from kinetic data | steps: rate law constants; initial rate}, group=TOAS07-P2-CHEM-S3-G3-4]
What is the initial rate of reaction (in M/min.) ?
\end{bankq}
\begin{bankq}[id=TOAS07-P2-CHEM-S4-Q07, source={TOAS - Test 07 (01 May 2022)}, paper=2, qno=7, type=numerical, subtopic={Arrhenius equation}, answer={3}, answer_src=key, difficulty=4, diff_parts={K2 S1 I1 C0 F1}, diff_note={concepts: Arrhenius equation, activation energy and enthalpy | steps: Ea forward = 11000 R = 21.9 kcal; Ea reverse = Ea forward - delta H | I: relation between the two activation energies}]
The rate constant for decomposition of $\mathrm{COCl_2(g)}$ according to reaction :\\
$\mathrm{COCl_2(g)\rightleftharpoons CO(g)+Cl_2(g)}$ ; $\Delta H=19$ kcal mole$^{-1}$\\
is expressed as : $\ln k=15-\dfrac{11000}{T(K)}$\\
Calculate activation energy for formation of $\mathrm{COCl_2}$ from $\mathrm{CO(g)}$ \& $\mathrm{Cl_2(g)}$ in kcal mole$^{-1}$ ?
\end{bankq}
\begin{bankq}[id=TOAS08-P1-CHEM-S3-Q03, source={TOAS - Test 08 (15 May 2022)}, paper=1, qno=3, type=numerical, subtopic={Order}, answer={1.33}, answer_src=key, difficulty=2, diff_parts={K1 S1 I0 C0 F1}, diff_note={concepts: rate of reaction from a graph | steps: rate of disappearance of A from the graph; divide by 4}]
For the reaction , $4\mathrm A\to3\mathrm B$, the following graph is obtained experimentally, the rate of reaction (in M/sec) when $[\mathrm A]=2$M is
\bfig{103.8}{TOAS08_p1_chem_s3_q03_1.png}
\end{bankq}
\begin{bankq}[id=TOAS08-P1-CHEM-S3-Q04, source={TOAS - Test 08 (15 May 2022)}, paper=1, qno=4, type=numerical, subtopic={Order}, answer={382.00 to 384.00}, answer_src=key, difficulty=4, diff_parts={K1 S2 I0 C1 F1}, diff_note={concepts: first order decomposition with total pressure | steps: initial pressure from the final pressure; rate constant from the 30 min value; pressure at 75 min including the vapour pressure of S}]
The decomposition of compound $P$, at temperature $T$ according to the equation $2P_{(g)}\to4Q_{(g)}+R_{(g)}+S_{(\ell)}$ is the first order reaction. After 30 minutes from the start of decomposition in a closed vessel, the total pressure developed is found to be 317 mm Hg and after a long period of time the total pressure observed to be 617 mm Hg. The total pressure of the vessel after 75 minute, if the volume of liquid $S$ is supposed to be negligible is: (Given : Vapour pressure of $S(\ell)$ at temperature $T=32.5$ mm Hg, $\ln(1.169)=0.16$, $\ln(1.5)=0.40$)
\end{bankq}
\begin{bankq}[id=TOAS08-P2-CHEM-S2-Q09, source={TOAS - Test 08 (15 May 2022)}, paper=2, qno=9, type=single, subtopic={Order}, answer={A}, answer_src=key, difficulty=3, diff_parts={K1 S2 I0 C1 F0}, diff_note={concepts: second order gas phase reaction | steps: initial concentration and rate constant; concentration and pressure at times; match}]
For the reaction $2\mathrm A(g)\to\mathrm B(g)$, starting with 4 moles of $A$ \& half life 40 sec in a 5 litre container at 400 K.
\begin{center}\small\renewcommand{\arraystretch}{1.1}\setlength{\extrarowheight}{3.5pt}\setlength{\tabcolsep}{3pt}
\begin{tabular}{|c|l|c|c|}
\hline
 & \multicolumn{1}{c|}{List-I} & & List-II\\[3pt]
\hline
(P) & \parbox[t]{4cm}{If reaction follow first order then initial rate of reaction} & (1) & $3.46\times10^{-4}$ M/s\\[3pt]
(Q) & \parbox[t]{4cm}{If reaction follow zero order then rate of reaction at time : $t=40$ sec.} & (2) & $6.93\times10^{-3}$ M/s\\[3pt]
(R) & \parbox[t]{4cm}{If reaction follow first order kinetics then rate of reaction at 95\% completion.} & (3) & $5\times10^{-3}$ M/s\\[3pt]
(S) & \parbox[t]{4cm}{If reaction follow second order kinetics then rate of reaction at 99\% completion} & (4) & $10^{-6}$ M/s\\[3pt]
 & & (5) & $6.25\times10^{-5}$ M/s\\[3pt]
\hline
\end{tabular}
\end{center}
\begin{qoptions}
\opt{A}{$P\to2;Q\to3;R\to1;S\to4$}
\opt{B}{$P\to1;Q\to2;R\to3;S\to4$}
\opt{C}{$P\to5;Q\to1;R\to2;S\to4$}
\opt{D}{$P\to2;Q\to1;R\to4;S\to3$}
\end{qoptions}
\end{bankq}
\begin{bankpassage}[id=TOAS09-P1-CHEM-S3-G3-4, source={TOAS - Test 09 (30 May 2022)}, questions={3-4}]
Kinetics of acidic hydrolysis of ester is pseudo I order reaction
\[\begin{aligned}&\mathrm{CH_3COOCH_3(aq.)+\underset{excess}{H_2O}\,(\ell)}\\&\qquad\xrightarrow{[H^+]}\mathrm{CH_3COOH(aq.)+CH_3OH(aq.)}\end{aligned}\]
Rate law is given by : $r=K[\mathrm{H^+}][\mathrm{H_2O}][\mathrm{CH_3COOCH_3}]$ ; $K=1.8\times10^{-3}\ \mathrm{M^{-2}\,sec^{-1}}$ where H$^+$ ion concentration is given by acid catalyst. In an experiment 1M CH$_3$COOCH$_3$ is hydrolysed using 0.1M HCl as catalyst. [Use molarity of H$_2$O = $\dfrac{1000}{18}$M]
\end{bankpassage}
\begin{bankq}[id=TOAS09-P1-CHEM-S3-Q03, source={TOAS - Test 09 (30 May 2022)}, paper=1, qno=3, type=numerical, subtopic={Order}, answer={91.00}, answer_src=key, difficulty=3, diff_parts={K1 S2 I0 C0 F1}, diff_note={concepts: first order kinetics by titration | steps: titre difference proportional to product formed; rate constant from the passage; titre at the given time}, group=TOAS09-P1-CHEM-S3-G3-4]
Reaction mixture is titrated against standard NaOH solution at different times. If initially 10 ml of NaOH was required while after long time 100 ml of NaOH is required then find volume of NaOH (in ml) required at 230.3 sec.
\end{bankq}
\begin{bankq}[id=TOAS09-P1-CHEM-S3-Q04, source={TOAS - Test 09 (30 May 2022)}, paper=1, qno=4, type=numerical, subtopic={Order}, answer={140.00}, answer_src=key, difficulty=2, diff_parts={K1 S1 I0 C0 F1}, diff_note={concepts: first order kinetics | steps: rate constant from half life; time to reach the concentration}, group=TOAS09-P1-CHEM-S3-G3-4]
Calculate the time at which ester concentration reduce to 0.25M (in seconds) [ln2 = 0.7]
\end{bankq}
\begin{bankq}[id=TOAS09-P2-CHEM-S1-Q17, source={TOAS - Test 09 (30 May 2022)}, paper=2, qno=17, type=single, subtopic={Order}, answer={B}, answer_src=key, difficulty=2, diff_parts={K1 S2 I0 C0 F0}, diff_note={concepts: second order kinetics | steps: rate law for second order in A; concentration at 50 percent conversion; rate}]
For a process A + B $\to$ P, the reaction is second order with respect to A and zero order with respect to B. When 1 mole each of A \& B are taken in 1 litre vessel the initial rate is $1\times10^{-2}$ mol/L-sec. The rate of reaction when 50\% of the reactant have been converted to product would be -
\begin{qoptions}
\opt{A}{$1\times10^{-2}$ mol/L-sec}
\opt{B}{$2.5\times10^{-3}$ mol/L-sec}
\opt{C}{$5\times10^{-2}$ mol/L-sec}
\opt{D}{$5\times10^{-3}$ mol/L-sec}
\end{qoptions}
\end{bankq}
\begin{bankq}[id=TOAS09-P2-CHEM-S2-Q09, source={TOAS - Test 09 (30 May 2022)}, paper=2, qno=9, type=numerical, subtopic={Rate laws}, answer={0.50}, answer_src=solved, working={$-d[A]/dt=x r$, $d[B]/dt=y r$; $\log(xr)=\log(yr)+\log2$ gives $x=2y$, $y/x=0.5$ (the printed key shows 1)}, difficulty=2, diff_parts={K1 S1 I0 C0 F1}, diff_note={concepts: stoichiometry of rates | steps: rate of disappearance and appearance in terms of x and y; log relation gives x = 2y}]
In the following reaction : $x\mathrm{A}\longrightarrow y\mathrm{B}$, following equation is given :
\[\log\left[-\frac{d[\mathrm{A}]}{dt}\right]=\log\left[\frac{d[\mathrm{B}]}{dt}\right]+0.3,\]
Find ratio $\left(\dfrac yx\right)$ : [Given $\log2=0.3$]
\end{bankq}
\begin{bankpassage}[id=TOAS10-P1-CHEM-S2-G9-10, source={TOAS - Test 10 (12 Jun 2022)}, questions={9-10}]
A substance `X' converts into `Z' through some intermediate, `Y'.
\bfig{193.7}{TOAS10_p1_chem_s2_para9_10_1.png}
\end{bankpassage}
\begin{bankq}[id=TOAS10-P1-CHEM-S2-Q09, source={TOAS - Test 10 (12 Jun 2022)}, paper=1, qno=9, type=multi, subtopic={Rate laws}, answer={A,B,C,D}, answer_src=key, difficulty=4, diff_parts={K1 S2 I0 C1 F1}, diff_note={concepts: consecutive first order reactions | steps: rate constants from the figure; [X] at 20 min; [Y] at 20 min; [Z] at 20 min and the sum}, group=TOAS10-P1-CHEM-S2-G9-10]
The correct concentration(s) of `X', `Y' and `Z' at $t=20$ min is ($e^{-1}=0.36$, $e^{-10}=4.5\times10^{-5}$) :
\begin{qoptions}
\opt{A}{[X] = 0.18 M}
\opt{B}{[Y] $\approx$ 0.02 M}
\opt{C}{[Z] $\approx$ 0.30 M}
\opt{D}{[X] + [Y] + [Z] = 0.50 M}
\end{qoptions}
\end{bankq}
\begin{bankq}[id=TOAS10-P1-CHEM-S2-Q10, source={TOAS - Test 10 (12 Jun 2022)}, paper=1, qno=10, type=multi, subtopic={Rate laws}, answer={A,B,C}, answer_src=key, difficulty=4, diff_parts={K1 S2 I1 C0 F1}, diff_note={concepts: maximum intermediate concentration in consecutive reactions | steps: t at maximum Y = ln(k1/k2)/(k1 - k2); rate of formation of Z; condition for equal rates; [X] value | I: maximum of Y occurs when dY/dt = 0}, group=TOAS10-P1-CHEM-S2-G9-10]
The correct information(s) is/are ($\ell$n 2 = 0.7, $\ell$n = 1.6) :
\begin{qoptions}
\opt{A}{The concentration of `Y' becomes maximum at about 5.11 min. from the start of reaction}
\opt{B}{The rate of formation of `Z' becomes maximum at about 5.11 min. from the start of reaction}
\opt{C}{The rate of disappearance of `X' becomes equal to the rate of appearance of `Z' at about 5.11 min. from the start of reaction}
\opt{D}{Concentration of `X' becomes 0.25 M at about 5.11 min. from the start of reaction}
\end{qoptions}
\end{bankq}
\begin{bankq}[id=TOAS11-P1-CHEM-S1-Q01, source={TOAS - Test 11 (27 Jun 2022)}, paper=1, qno=1, type=single, subtopic={Order}, answer={A}, answer_src=key, difficulty=2, diff_parts={K1 S2 I0 C0 F0}, diff_note={concepts: first order gaseous decomposition | steps: pressure of A from the total pressure at 30 min; rate constant; time for 50 mm}]
For a first order reaction A $\to$ B + C all the reactant and product molecules are in gaseous state. If the initial pressure of the reactant molecules is 100 mm and after 30 minutes the total pressure of the mixture containing molecules of all three A, B \& C gases is 140 mm, then find the time interval at which initial pressure of A will be reduced to 50 mm. ($\log_{10}3=0.5$, $\log_{10}2=0.3$)
\begin{qoptions}
\opt{A}{45 min}
\opt{B}{30 min}
\opt{C}{60 min}
\opt{D}{15 min}
\end{qoptions}
\end{bankq}
\begin{bankq}[id=TOAS11-P1-CHEM-S1-Q02, source={TOAS - Test 11 (27 Jun 2022)}, paper=1, qno=2, type=single, subtopic={Rate laws}, answer={D}, answer_src=key, difficulty=2, diff_parts={K1 S2 I0 C0 F0}, diff_note={concepts: rate law from a mechanism | steps: rate determining step; equilibrium for N2O2; expression for each rate and K}]
For the reaction:\par
2H$_2$ + 2NO $\rightleftharpoons$ N$_2$ + 2H$_2$O\par
the following mechanism has been suggested :\par
2NO $\rightleftharpoons$ N$_2$O$_2$ [Equilibrium Constant = K$_1$ (fast)]\par
N$_2$O$_2$ + H$_2$ $\xrightarrow{k_2}$ N$_2$O + H$_2$O (Slow)\par
N$_2$O + H$_2$ $\xrightarrow{k_3}$ N$_2$ + H$_2$O (fast)\par
According to above mechanism which option is correctly matched:
\begin{qoptions}
\opt{A}{$-\dfrac12\dfrac{d[\mathrm{H_2}]}{dt}=K[\mathrm{NO}][\mathrm{H_2}]$\quad Where $K=k_2\times k_3$}
\opt{B}{$-\dfrac12\dfrac{d[\mathrm{NO}]}{dt}=K[\mathrm{NO}][\mathrm{H_2}]$\quad Where $K=k_2\times K_1$}
\opt{C}{$\dfrac{d[\mathrm{N_2}]}{dt}=K[\mathrm{NO}][\mathrm{H_2}]$\quad Where $K=k_2\times K_1$}
\opt{D}{$-\dfrac12\dfrac{d[\mathrm{NO}]}{dt}=K[\mathrm{NO}]^2[\mathrm{H_2}]$\quad Where $K=k_2\times K_1$}
\end{qoptions}
\end{bankq}
\begin{bankq}[id=TOAS11-P2-CHEM-S3-Q01, source={TOAS - Test 11 (27 Jun 2022)}, paper=2, qno=1, type=single, subtopic={Arrhenius equation}, answer={13.02}, answer_src=key, difficulty=1, diff_parts={K1 S1 I0 C0 F0}, diff_note={concepts: Arrhenius equation | steps: ln(k2/k1) = Ea/R (1/T1 - 1/T2); solve Ea}]
If rate of reaction becomes double when temperature is increased from 27$^\circ$C to 37$^\circ$C. Then find activation energy of reaction in kcal/mole. ($R=2$ cal/mol-K) ($\ln2=0.7$)
\end{bankq}
\begin{bankq}[id=TOAS12-P2-CHEM-S1-Q02, source={TOAS - Test 12 (10 Jul 2022)}, paper=2, qno=2, type=single, subtopic={Order}, answer={A}, answer_src=key, difficulty=2, diff_parts={K1 S2 I0 C0 F0}, diff_note={concepts: order from pressure change | steps: pressure rise for each order; time condition; which orders are possible}]
Consider the following reaction $2A(g)\to3B(g)+C(g)$. Starting with pure A having pressure 2 atm initially, the total pressure is exactly doubled in 2 hrs. The possible order of reaction is :-
\begin{qoptions}
\opt{A}{zero}
\opt{B}{first}
\opt{C}{second}
\opt{D}{third}
\end{qoptions}
\end{bankq}
\begin{bankq}[id=TOAS13-P1-CHEM-S1-Q02, source={TOAS - Test 13 (24 Jul 2022)}, paper=1, qno=2, type=multi, subtopic={Rate laws}, answer={A,C,D}, answer_src=key, difficulty=3, diff_parts={K1 S2 I0 C0 F1}, diff_note={concepts: integrated rate law of order -1 | steps: integrate -dA/dt = k/A to A\^{}2 = A0\^{}2 - 2kt; half-life 3 A0\^{}2/8k; completion time A0\^{}2/2k; negative order means a complex reaction}]
For the reaction : A $\to$ B. Differential rate law is : $\dfrac{-d[A]}{dt}=k[A]^{-1}$. If initial concentration of A is $A_0$M and concentration at time `t' is $A_t$ M then identify correct option (s)
\begin{qoptions}
\opt{A}{$A_t^2=A_0^2-2kt$}
\opt{B}{Half life, $t_{1/2}=\dfrac{3A_0^2}{4K}$}
\opt{C}{Completion time for reaction : $t_{100\%}=\dfrac{A_0^2}{2K}$}
\opt{D}{The reaction must be a complex reaction}
\end{qoptions}
\end{bankq}
\begin{bankpassage}[id=TOAS13-P1-CHEM-S2-G9-10, source={TOAS - Test 13 (24 Jul 2022)}, questions={9-10}]
Consider a three step mechanism\par
\noindent(i) I$_2$ $\underset{k_{-1}}{\overset{k_1}{\rightleftharpoons}}$ 2I \quad (fast)\par
\noindent(ii) H$_2$ $+$ I $\underset{k_{-2}}{\overset{k_2}{\rightleftharpoons}}$ H$_2$I \quad (fast)\par
\noindent(iii) H$_2$I $+$ I $\xrightarrow{k_3}$ 2HI \quad (slowest)
\end{bankpassage}
\begin{bankq}[id=TOAS13-P1-CHEM-S2-Q09, source={TOAS - Test 13 (24 Jul 2022)}, paper=1, qno=9, type=multi, subtopic={Rate laws}, answer={A,B,C}, answer_src=key, difficulty=5, diff_parts={K2 S2 I1 C0 F1}, diff_note={concepts: rate law from pre-equilibria, molecularity | steps: rate from the slow step; eliminate H2I and I using the fast equilibria; overall order 2; molecularity not defined for the overall reaction}, group=TOAS13-P1-CHEM-S2-G9-10]
Which is/are incorrect regarding above mechanism ?
\begin{qoptions}
\opt{A}{Rate of reaction $=K[\text{H}_2\text{I}]\,[\text{I}]$ where $K=k_3\cdot\dfrac{k_2}{k_{-2}}\cdot\left(\dfrac{k_1}{k_{-1}}\right)$}
\opt{B}{Rate is decided only by the steps (i) and (ii)}
\opt{C}{Molecularity of overall reaction is 2}
\opt{D}{Order of overall reaction is 2.}
\end{qoptions}
\end{bankq}
\begin{bankq}[id=TOAS13-P1-CHEM-S2-Q10, source={TOAS - Test 13 (24 Jul 2022)}, paper=1, qno=10, type=multi, subtopic={Rate laws}, answer={A,B,D}, answer_src=key, difficulty=5, diff_parts={K2 S2 I1 C0 F1}, diff_note={concepts: activation energy ordering, reaction profile | steps: k1 >> k2 >> k3 means the barriers increase; overall exothermic; identify the incorrect profiles}, group=TOAS13-P1-CHEM-S2-G9-10]
Which of the following qualitative plot of potential energy (PE) v/s reaction coordinates is/are incorrect if $k_1>>k_2>>k_3$ and reaction is exothermic :
\begin{qoptions}
\opt{A}{\bfigi{97.7}{TOAS13_p1_chem_s2_q10_A.png}}
\opt{B}{\bfigi{100.4}{TOAS13_p1_chem_s2_q10_B.png}}
\opt{C}{\bfigi{113.6}{TOAS13_p1_chem_s2_q10_C.png}}
\opt{D}{\bfigi{108.3}{TOAS13_p1_chem_s2_q10_D.png}}
\end{qoptions}
\end{bankq}
\begin{bankq}[id=TOAS13-P2-CHEM-S1-Q03, source={TOAS - Test 13 (24 Jul 2022)}, paper=2, qno=3, type=single, subtopic={Rate laws}, answer={C}, answer_src=key, difficulty=1, diff_parts={K1 S1 I0 C0 F0}, diff_note={concepts: reaction profile for a slow first step | steps: first barrier is higher; exothermic so products lie lower}]
For an exothermic reaction whose mechanism is given as :\par
A $+$ B $\to$ X \ (slow)\par
X $\to$ AB \ (fast)\par
The progress of the reaction can be best described by :
\begin{qoptions}
\opt{A}{\bfigi{63.2}{TOAS13_p2_chem_s1_q03_A.png}}
\opt{B}{\bfigi{63.4}{TOAS13_p2_chem_s1_q03_B.png}}
\opt{C}{\bfigi{63.4}{TOAS13_p2_chem_s1_q03_C.png}}
\opt{D}{\bfigi{63.4}{TOAS13_p2_chem_s1_q03_D.png}}
\end{qoptions}
\end{bankq}
\begin{bankq}[id=TOAS13-P2-CHEM-S1-Q18, source={TOAS - Test 13 (24 Jul 2022)}, paper=2, qno=18, type=single, subtopic={Order}, answer={D}, answer_src=key, difficulty=2, diff_parts={K1 S2 I0 C0 F0}, diff_note={concepts: first-order kinetics with pressure | steps: 75 minutes is three half-lives so A is 1/8; products from stoichiometry; total pressure 4.5 atm}]
For a first order gaseous reaction half life is 25 minutes and reaction started with 1 atm pressure of pure A, then find final pressure of reaction mixture after 75 minutes.\par
A(g) $\to$ 4B(g) $+$ C(g)
\begin{qoptions}
\opt{A}{2 atm}
\opt{B}{3 atm}
\opt{C}{3.75 atm}
\opt{D}{4.5 atm}
\end{qoptions}
\end{bankq}
\begin{bankq}[id=TOAS15-P1-CHEM-S2-Q08, source={TOAS - Test 15 (22 Aug 2022)}, paper=1, qno=8, type=multi, subtopic={Rate laws}, answer={B,C}, answer_src=key, difficulty=4, diff_parts={K2 S2 I0 C0 F1}, diff_note={concepts: catalyst and molecularity and order, Arrhenius equation | steps: catalyst does not change the equilibrium yield; elementary steps have order equal to molecularity; catalyst lowers the barrier; k increases with T}]
Select the correct statements
\begin{qoptions}
\opt{A}{Catalyst can change the percentage yield of a reversible reaction}
\opt{B}{Molecularity and order are same for the elementary reactions}
\opt{C}{The catalyst forms an activated complex of lower potential energy, with the reactants by which more number of molecules are able to cross the barrier per unit of time.}
\opt{D}{For exothermic reaction rate constant decreases with increasing temperature}
\end{qoptions}
\end{bankq}
\begin{bankq}[id=TOAS15-P2-CHEM-S2-Q06, source={TOAS - Test 15 (22 Aug 2022)}, paper=2, qno=6, type=numerical, subtopic={Rate laws}, answer={4.00}, answer_src=key, difficulty=2, diff_parts={K1 S1 I0 C0 F1}, diff_note={concepts: rate law of a first-order decomposition | steps: convert min\^{}-1 to s\^{}-1; rate = k C}]
The rate constant of the reaction $2\text{H}_2\text{O}_2\text{(aq)}\to2\text{H}_2\text{O(l)}+\text{O}_2\text{(g)}$ is $3\times10^{-3}$ min$^{-1}$. At what concentration of H$_2$O$_2$ (in M) the rate of the reaction will be $2\times10^{-4}$ M s$^{-1}$.
\end{bankq}
\begin{bankq}[id=TOAS15-P2-CHEM-S3-Q02, source={TOAS - Test 15 (22 Aug 2022)}, paper=2, qno=2, type=numerical, subtopic={Order}, answer={7}, answer_src=key, difficulty=5, diff_parts={K2 S2 I0 C1 F1}, diff_note={concepts: half-lives of first, zero order reactions | steps: half-life of the first order; half-life of the zero order; ratio of the constants}]
Consider two reactions at $27^\circ$C\par
(I). A $\to$ B ; ( I order kinetics)\par
(II). X $\to$ Y; ( Zero order kinetics)\par
If both A \& X each having initial concentration 0.1M \& same initial half life period. If ratio of rate constants for I \& II reactions $\left(\dfrac{K_I}{K_{II}}\right)$ is `X' find value of $\left(\dfrac X2\right)$ ? [$\ln2=0.7$]
\end{bankq}
\begin{bankq}[id=TOAS17-P2-CHEM-S1-Q05, source={TOAS - Test 17 (18 Sep 2022)}, paper=2, qno=5, type=single, subtopic={Order}, answer={A}, answer_src=key, difficulty=2, diff_parts={K1 S1 I1 C0 F0}, diff_note={concepts: order from a half-life plot | steps: slope of the log t against log a plot gives 1 - n; intercept gives k}]
What will be the order of reaction [A $\to$ P] and rate constant for this chemical change having [$\log t_{50\%}$] vs [$\log$ initial concentration of A (a)] curves as :
\bfig{112.3}{TOAS17_p2_chem_s1_q05_1.png}
\begin{qoptions}
\opt{A}{0, 1/2}
\opt{B}{1, 1}
\opt{C}{2, 2}
\opt{D}{3, 1}
\end{qoptions}
\end{bankq}
\begin{bankq}[id=TOAS17-P2-CHEM-S1-Q20, source={TOAS - Test 17 (18 Sep 2022)}, paper=2, qno=20, type=single, subtopic={Rate laws}, answer={A}, answer_src=key, difficulty=3, diff_parts={K1 S2 I1 C0 F0}, diff_note={concepts: acid-catalysed hydrolysis | steps: rate proportional to hydrogen ion concentration; weak acid ionisation; Ka}]
The initial rate of hydrolysis of methyl acetate (1M) by a weak acid (HA, 1M) is $1/100^{\text{th}}$ of that of a strong acid (HX, 1M), at 25$^\circ$C. The $K_a$ of HA is
\begin{qoptions}
\opt{A}{$1\times10^{-4}$}
\opt{B}{$1\times10^{-5}$}
\opt{C}{$1\times10^{-6}$}
\opt{D}{$1\times10^{-3}$}
\end{qoptions}
\end{bankq}
\begin{bankq}[id=TOAS18-P1-CHEM-S3-Q04, source={TOAS - Test 18 (30 Oct 2022)}, paper=1, qno=4, type=numerical, subtopic={Order}, answer={0.25}, answer_src=solved, working={Zero order: $t_{1/2}=C_0/2k$ gives $k=1$ mol L$^{-1}$ h$^{-1}$; time for 0.50 to 0.25 is $0.25/1=0.25$ h}, difficulty=2, diff_parts={K1 S1 I0 C0 F1}, diff_note={concepts: zero order kinetics | steps: rate constant from the half-life; time for the concentration change}]
The time for half-life period of a certain reaction. A $\to$ products is 1 h, when the initial concentration of the reactant `A', is 2.0 mol L$^{-1}$. How much time does it take (in hours) for its concentration to come from 0.50 to 0.25 mol L$^{-1}$, if it is a zero order reaction?
\end{bankq}
\begin{bankq}[id=TOAS18-P2-CHEM-S3-Q03, source={TOAS - Test 18 (30 Oct 2022)}, paper=2, qno=3, type=single, subtopic={Order}, answer={0}, answer_src=key, difficulty=3, diff_parts={K1 S2 I1 C0 F0}, diff_note={concepts: order from concentration-time data | steps: rate in each interval; compare; order}]
The concentration of R in the reaction R $\to$ P was measured as a function of time and the following data is obtained :
\begin{center}\footnotesize\setlength{\tabcolsep}{3pt}
\begin{tabular}{|l|c|c|c|c|}\hline
[R] (molar)&1.0&0.75&0.40&0.10\\\hline
$t$ (min.)&0.0&0.05&0.12&0.18\\\hline
\end{tabular}\end{center}
The order of the reaction is
\end{bankq}

\section*{TAAS papers}
\begin{bankq}[id=TAAS01-P1-CHEM-S2-Q07, source={TAAS - Test 01 (18 Jun 2023)}, paper=1, qno=7, type=multi, subtopic={Order}, answer={A,B}, answer_src=key, difficulty=5, diff_parts={K1 S2 I1 C1 F1}, diff_note={concepts: stoichiometry and rate of a gas phase reaction | steps: pressures from stoichiometry; progress at each time; order; end of reaction}, also_in={TNAS - Test 04 (18 Jun 2023) P1 Q5}]
$\mathrm{2A(g)\rightarrow4B(g)+C(g)}$

At $t=0$ pressure of A is 2 atm, at $t=10$ sec. pressure of B is 2 atm and at $t=30$ sec. pressure of C is 0.75 atm. Select correct option/options :
\begin{qoptions}
\opt{A}{At $t=70$ sec pressure of B will be 3.5 atm.}
\opt{B}{Reaction may be an elementary reaction}
\opt{C}{At $t=70$ sec, rate of formation of B is $\dfrac1{80}$ atm/sec.}
\opt{D}{Reaction will complete in definite time.}
\end{qoptions}
\end{bankq}
\begin{bankq}[id=TAAS01-P1-CHEM-S3-Q04, source={TAAS - Test 01 (18 Jun 2023)}, paper=1, qno=4, type=single, subtopic={Order}, answer={120.00}, answer_src=key, difficulty=2, diff_parts={K1 S2 I0 C0 F0}, diff_note={concepts: second-order kinetics | steps: rate constant from the half-life; integrated second-order rate law; time}]
The decomposition of A$_2$B$_2$ follows the rate equation $r=K[\mathrm{A_2B_2}]^2$.

$\mathrm{A_2B_2}\xrightarrow{K}\mathrm{A_2+B_2}$. The kinetic data of the reaction is given below.
\begin{center}\footnotesize\setlength{\tabcolsep}{3pt}
\begin{tabular}{|L{3.3cm}|L{1.1cm}|}\hline
Initial conc. of A$_2$B$_2$ (mole/lit)&7\\\hline
First half life of A$_2$B$_2$ in sec.&1440\\\hline
\end{tabular}\end{center}
If the initial concentration of A$_2$B$_2$ were 56 millimol/lit, Calculate the time (in sec) taken for the conversion of A$_2$B$_2$ to reach 60\% of initial value.

(Report your answer as $\dfrac t{1000}$ )
\end{bankq}
\begin{bankq}[id=TAAS02-P1-CHEM-S2-Q05, source={TAAS - Test 02 (30 Jul 2023)}, paper=1, qno=5, type=multi, subtopic={Order}, answer={A,C}, answer_src=key, difficulty=3, diff_parts={K1 S2 I0 C0 F1}, diff_note={concepts: parallel first-order reactions | steps: ratio of products; effect of a catalyst; overall rate constant and activation energy}]
Consider the following parallel, first order processes
\bfigi{41.8}{TAAS02_p1_chem_s2_q05_1.png} $k_1$, $k_2$ are rate constants and $E_1$ and $E_2$ are activation energies.

The correct deduction(s) is/are :
\begin{qoptions}
\opt{A}{At any time, the ratio of concentration of P and Q has a fixed value at a given temperature.}
\opt{B}{If a catalyst is added, it will always decrease both $E_1$ and $E_2$ but by same amount.}
\opt{C}{If `$E$' represents the overall activation energy for decay of A, it will be $\left(\dfrac{E_1k_1+E_2k_2}{k_1+k_2}\right)$.}
\opt{D}{If `$k$' represents the overall rate constant for decay of A, it will be $\dfrac1k=\dfrac1{k_1}+\dfrac1{k_2}$.}
\end{qoptions}
\end{bankq}
\begin{bankq}[id=TAAS02-P2-CHEM-S1-Q04, source={TAAS - Test 02 (30 Jul 2023)}, paper=2, qno=4, type=single, subtopic={Arrhenius equation}, answer={D}, answer_src=key, difficulty=1, diff_parts={K1 S1 I0 C0 F0}, diff_note={concepts: catalysis | steps: catalysts lower activation energy; do not change the enthalpy}]
Which of the following explains the increase of reaction rate by a catalyst ?
\begin{qoptions}
\opt{A}{Catalyst provides the necessary energy to the colliding molecules to cross the barrier.}
\opt{B}{Catalyst decreases the rate of backward reaction so that the rate of forward reaction increases.}
\opt{C}{Catalyst decreases the enthalpy change of the reaction.}
\opt{D}{Catalyst provides an alternative path of lower activation energy.}
\end{qoptions}
\end{bankq}

\section*{TEAS papers}
\begin{bankq}[id=TEAS01-P1-CHEM-S3-Q01, source={TEAS - Test 01 (16 Apr 2023)}, paper=1, qno=1, type=numerical, subtopic={Rate laws}, answer={1.00}, answer_src=key, difficulty=5, diff_parts={K1 S2 I1 C1 F1}, diff_note={concepts: rate constant from titration | steps: concentration proportional to the difference of volumes; first order rate law; pseudo constant; K}]
In the acid catalysed hydrolysis of ethoxy ethanoate the rate law is found to be Rate $=K[\mathrm{CH_3COOC_2H_5}][\mathrm{H^+}]$. Calculate the value of ``$K$'' in $\mathrm{M^{-1}\,min^{-1}}$ from the following titration data.\par\smallskip
{\footnotesize\setlength{\tabcolsep}{3pt}\begin{tabular}{|p{3.0cm}|l|l|l|}
\hline
\multicolumn{4}{|l|}{$[\mathrm{H^+}]=0.1$ M}\\
\hline
Time (min) & $t=0$ & $t=6.93$ min & $t=\infty$\\
\hline
Volume of NaOH required for complete neutralization of small and equal volume of reaction mixture taken out at different time interval. & 35 ml & 60 ml & 85 ml\\
\hline
\end{tabular}}
\end{bankq}

\section*{TMAS papers}
\begin{bankq}[id=TMAS01-P1-CHEM-S1-Q12, source={TMAS - Test 01 (09 Jul 2023)}, paper=1, qno=12, type=single, subtopic={Rate laws}, answer={B}, answer_src=key, difficulty=4, diff_parts={K1 S2 I1 C1 F0}, diff_note={concepts: first order kinetics from titration | steps: volumes proportional to concentration; rate constant; half life}]
Methyl acetate was hydrolyzed in HCl aqueous solution at 298K, samples of reaction mixture were withdrawn at different time intervals and titrated with 0.12N NaOH sample whose volumes were obtained as :\par\smallskip
{\footnotesize\setlength{\tabcolsep}{4pt}\begin{tabular}{|l|c|c|c|}
\hline
Time(mins) & 0 & 25 & $\infty$\\
\hline
Volume of NaOH(mL) & 18 & 30 & 42\\
\hline
\end{tabular}}\par\smallskip
Half- life for the reaction below is :\\
\par\noindent\hfill{\footnotesize$\mathrm{CH_3COOCH_3(aq.)+H_2O(\ell)}\xrightarrow{\mathrm{H^\oplus}}\mathrm{CH_3COOH(aq.)+CH_3OH(aq.)}$}\hfill\null\par
\begin{qoptions}
\opt{A}{12.5 minutes}
\opt{B}{25 minutes}
\opt{C}{37.5 minutes}
\opt{D}{50 minutes}
\end{qoptions}
\end{bankq}
\begin{bankq}[id=TMAS01-P1-CHEM-S2-Q06, source={TMAS - Test 01 (09 Jul 2023)}, paper=1, qno=6, type=numerical, subtopic={Order}, answer={120}, answer_src=key, difficulty=6, diff_parts={K1 S2 I2 C1 F1}, diff_note={concepts: second order kinetics | steps: half life depends on concentration; rate constant; time for 40 percent; scaled}]
The decomposition of $\mathrm{A_2B_2}$ follows the rate equation $r=K[\mathrm{A_2B_2}]^2$.\\
$\mathrm{A_2B_2\xrightarrow{K}A_2+B_2}$. The kinetic data of the reaction is given below.\par\smallskip
{\footnotesize\setlength{\tabcolsep}{3pt}\begin{tabular}{|l|c|}
\hline
Initial conc. of $\mathrm{A_2B_2}$ (mole/lit) & 7\\
\hline
First half life of $\mathrm{A_2B_2}$ in sec. & 1440\\
\hline
\end{tabular}}\par\smallskip
If the initial concentration of $\mathrm{A_2B_2}$ were 56 m mol/lit, Calculate the time (in sec) taken for the conversion of $\mathrm{A_2B_2}$ to reach 60\% of initial value.\\
(Report your answer as $\dfrac t{1000}$ )
\end{bankq}
\begin{bankpassage}[id=TMAS01-P2-CHEM-S2-G8-10, source={TMAS - Test 01 (09 Jul 2023)}, questions={8-10}]
Initially `$a$' moles of $A$ react with 2 moles of $B$ in a 1 litre container according to given reactions for different values of `$a$' :

\noindent\begin{minipage}{\textwidth}
\footnotesize\setlength{\tabcolsep}{4pt}
\begin{tabular}{|l|p{5.2cm}|l|p{3.2cm}|l|p{4.6cm}|}
\hline
\multicolumn{2}{|c|}{Column-I} & \multicolumn{2}{c|}{Column-II} & \multicolumn{2}{c|}{Column-III}\\
\hline
(I) & $\mathrm{A(g)+2B(g)\to C(g)+D(g)}$\newline $r=(2\times10^{-3})[A]^{1/2}[B]^{1/2}$\newline $a=1$ & (i) & Has highest value of half life & (P) & May be an elementary reaction\\
\hline
(II) & $\mathrm{2A(g)+B(g)\to C(g)+2D(g)}$\newline $r=(2\times10^{-3})[A]^{1}[B]^{0}$\newline $a=4$ & (ii) & Has lowest value of half life & (Q) & Must be a complex reaction\\
\hline
(III) & $\mathrm{A(g)+B(g)\to 2C(g)+D(g)}$\newline $r=(2\times10^{-3})[A][B]$\newline $a=2$ & (iii) & Reaction is complete in definite time & (R) & After equal time intervals remaining concentrations of $B$ will be in geometric progression\\
\hline
(IV) & $\mathrm{A+3B\to 2C(g)+3D(g)}$\newline $r=(2\times10^{-3})[A]^{0}[B]^{2}$\newline $a=(2/3)$ & (iv) & Reaction is complete in infinite time & (S) & After equal time intervals remaining concentration of $B$ will be in harmonic progression\\
\hline
\end{tabular}
\end{minipage}
\end{bankpassage}
\begin{bankq}[id=TMAS01-P2-CHEM-S2-Q08, source={TMAS - Test 01 (09 Jul 2023)}, paper=2, qno=8, type=single, subtopic={Order}, answer={C}, answer_src=key, difficulty=5, diff_parts={K1 S2 I2 C1 F0}, diff_note={concepts: order and half life | steps: order of each; time to completion; elementary or complex; match}, group=TMAS01-P2-CHEM-S2-G8-10]
For given reactions only correct option will be :
\begin{qoptions}
\opt{A}{(I) (ii) (P)}
\opt{B}{(II) (iii) (R)}
\opt{C}{(III) (i) (P)}
\opt{D}{(IV) (ii) (P)}
\end{qoptions}
\end{bankq}
\begin{bankq}[id=TMAS01-P2-CHEM-S2-Q09, source={TMAS - Test 01 (09 Jul 2023)}, paper=2, qno=9, type=single, subtopic={Order}, answer={A}, answer_src=key, difficulty=5, diff_parts={K1 S2 I2 C1 F0}, diff_note={concepts: order and half life | steps: order of each; time to completion; elementary or complex; match}, group=TMAS01-P2-CHEM-S2-G8-10]
For given reactions only incorrect option will be :
\begin{qoptions}
\opt{A}{(I) (iii) (Q)}
\opt{B}{(II) (iv) (R)}
\opt{C}{(III) (iv) (P)}
\opt{D}{(IV) (ii) (Q)}
\end{qoptions}
\end{bankq}
\begin{bankq}[id=TMAS01-P2-CHEM-S2-Q10, source={TMAS - Test 01 (09 Jul 2023)}, paper=2, qno=10, type=single, subtopic={Order}, answer={B}, answer_src=key, difficulty=5, diff_parts={K1 S2 I2 C1 F0}, diff_note={concepts: order and half life | steps: order of each; time to completion; elementary or complex; match}, group=TMAS01-P2-CHEM-S2-G8-10]
For (I) correct option will be :
\begin{qoptions}
\opt{A}{(I) (iii) (Q)}
\opt{B}{(I) (iv) (R)}
\opt{C}{(I) (i) (S)}
\opt{D}{(I) (ii) (Q)}
\end{qoptions}
\end{bankq}
\begin{bankq}[id=TMAS02-P1-CHEM-S2-Q05, source={TMAS - Test 02 (06 Aug 2023)}, paper=1, qno=5, type=multi, subtopic={Rate laws}, answer={A,B,C,D}, answer_src=key, difficulty=6, diff_parts={K2 S2 I1 C1 F1}, diff_note={concepts: reaction mechanism, order | steps: intermediates; catalyst; rate determining step; order}]
For the following mechanism one of the involved species is a catalyst :\\
\textbf{Step (i) :} {\small$\mathrm{H^++HNO_2}\mathrel{\overset{K_1}{\underset{K_2}{\rightleftharpoons}}}\mathrm{H_2NO_2^+}$}\\
\textbf{Step (ii) :} {\small$\mathrm{H_2NO_2^++Br^-}\mathrel{\overset{K_3}{\rightleftharpoons}}\mathrm{NOBr+H_2O}$}\\
\textbf{Step (iii) :}\\ {\small$\mathrm{NOBr+C_6H_5NH_2}\mathrel{\overset{K_4}{\rightleftharpoons}}\mathrm{C_6H_5N_2^++Br^-+H_2O}$}\\
Select the correct option(s) :
\begin{qoptions}
\opt{A}{Intermediates are $\mathrm{H_2NO_2^+}$ and $\mathrm{NOBr}$}
\opt{B}{Catalyst is $\mathrm{Br^-}$}
\opt{C}{If step (ii) forward reaction is slowest step then overall order of reaction is $3$.}
\opt{D}{If step (i) forward reaction is slowest step then overall order of reaction is $2$.}
\end{qoptions}
\end{bankq}
\begin{bankq}[id=TMAS04-P1-CHEM-S3-Q02, source={TMAS - Test 04 (19 Nov 2023)}, paper=1, qno=2, type=numerical, subtopic={Arrhenius equation}, answer={10.00}, answer_src=key, difficulty=6, diff_parts={K1 S2 I2 C1 F1}, diff_note={concepts: Arrhenius plot | steps: read intercept and slope; convert; rate constant}]
Rate constant ($k$) of elementary reaction $\mathrm{A(g)}\to\mathrm{B(g)}$ at $27^\circ$C is \underline{\hspace{1cm}}$\times10^{7}$ (sec$^{-1}$). Following plot is given for this reaction with units of $k$ as sec$^{-1}$.
\bfig{109.0}{TMAS04_p1_chem_s3_q02_1.png}
\end{bankq}

\section*{TNAS papers}
\begin{bankq}[id=TNAS03-P1-CHEM-S1-Q07, source={TNAS - Test 03 (28 May 2023)}, paper=1, qno=7, type=single, subtopic={Order}, answer={A}, answer_src=key, difficulty=3, diff_parts={K1 S1 I1 C1 F0}, diff_note={concepts: enzyme kinetics graph | steps: rate against substrate concentration; orders at low and high concentration}]
The kinetics of an enzyme-catalyzed hydrolysis reaction is represented by the following graph, where $[S]$ is the substrate concentration, and $v$ is the rate of reaction :
\bfig{71.5}{TNAS03_p1_chem_s1_q07_1.png}
The kinetic course of the reaction can be described as :
\begin{qoptions}
\opt{A}{First order, zero order.}
\opt{B}{Zero order, first order}
\opt{C}{First order, second order}
\opt{D}{Second order, first order}
\end{qoptions}
\end{bankq}
\begin{bankq}[id=TNAS03-P1-CHEM-S1-Q16, source={TNAS - Test 03 (28 May 2023)}, paper=1, qno=16, type=single, subtopic={Arrhenius equation}, answer={B}, answer_src=key, difficulty=4, diff_parts={K1 S2 I1 C1 F0}, diff_note={concepts: reaction profile | steps: intermediates; transition states; slowest step}]
Consider the following reaction that goes from A to B in three steps as shown below :
\bfig{155.5}{TNAS03_p1_chem_s1_q16_1.png}
Choose the correct\\
\begin{center}
{\small\setlength{\tabcolsep}{3pt}\begin{tabular}{|p{0.22\linewidth}|p{0.25\linewidth}|p{0.22\linewidth}|}\hline
Number of intermediates & Number of Activated complexes & Rate determining step\\\hline
\end{tabular}}
\end{center}
\begin{qoptions}
\opt{A}{$3$ \quad $2$ \quad II}
\opt{B}{$2$ \quad $3$ \quad II}
\opt{C}{$2$ \quad $3$ \quad I}
\opt{D}{$2$ \quad $3$ \quad III}
\end{qoptions}
\end{bankq}
\begin{bankq}[id=TNAS03-P1-CHEM-S2-Q08, source={TNAS - Test 03 (28 May 2023)}, paper=1, qno=8, type=numerical, subtopic={Order}, answer={600}, answer_src=key, difficulty=6, diff_parts={K1 S2 I2 C1 F1}, diff_note={concepts: fractional order kinetics | steps: rate law form; integrate; time to reach concentration}]
Reaction $\mathrm{A+B\to C+D}$ follow's following rate law : rate $=k=[A]^{+\frac12}[B]^{\frac12}$. Starting with initial conc. of one mole of A and B each, what is the time taken for amount of A to become 0.25 mole. [Given $k=2.31\times10^{-3}$ sec$^{-1}$]\\
($\ln2=0.693$)
\end{bankq}
\begin{bankq}[id=TNAS03-P2-CHEM-S3-Q08, source={TNAS - Test 03 (28 May 2023)}, paper=2, qno=8, type=numerical, subtopic={Order}, answer={4.75}, answer_src=key, difficulty=6, diff_parts={K1 S2 I2 C1 F1}, diff_note={concepts: kinetics with stoichiometric ratio | steps: first order in B; rate constant from half life; moles reacted; concentration of A}]
The reaction $\mathrm{2A+B_2\to2\,AB}$ is first order with respect to $[B_2]$ and zero order with respect to $[A]$. The half-life for the overall reaction is 2 minutes when concentrations of $[A]$ and $[B]$ are in the 2 : 1 molar ratio. If 0.100 moles of A and 0.0300 moles of $\mathrm B_2$ are dissolved in 10 mL of solvent, what will concentration of $[A]$ be, 6.00 minutes after the reaction begins ?
\end{bankq}
\begin{bankq}[id=TNAS04-P1-CHEM-S2-Q07, source={TNAS - Test 04 (18 Jun 2023)}, paper=1, qno=7, type=multi, subtopic={Order}, answer={C,D}, answer_src=key, difficulty=8, diff_parts={K3 S3 I1 C1 F1}, diff_note={concepts: activated complex, half life, consecutive and parallel first order reactions | steps: statement on activated complex; pseudo order argument; maximum of the intermediate; rate formation maximum; parallel ratio}]
Choose correct statement(s) :
\begin{qoptions}
\opt{A}{All activated complex get converted into products.}
\opt{B}{Half life of `A', would be independent of initial concentration of `B' in an elementary reaction : $\mathrm{A(g)+2B(g)\longrightarrow}$ Products which starts with 1 mol `A' and 1000 moles of B in 1 litre vessel.}
\opt{C}{In I order reaction : $\mathrm{A\xrightarrow{K_1}B\xrightarrow{K_2}C}$ rate of formation of C will be maximum at time $\dfrac{\ln(K_2/K_1)}{(K_2-K_1)}$}
\opt{D}{A parallel first order reaction \bfigi{47.5}{TNAS04_p1_chem_s2_q07_1.png} $K_1:K_2=\dfrac21$. Mole ratio of `B' \& `C' would be $3:2$, when 50\% of `A' has reacted.}
\end{qoptions}
\end{bankq}
\begin{bankq}[id=TNAS04-P2-CHEM-S1-Q01, source={TNAS - Test 04 (18 Jun 2023)}, paper=2, qno=1, type=multi, subtopic={Arrhenius equation}, answer={A,B,C}, answer_src=key, difficulty=6, diff_parts={K2 S2 I1 C1 F1}, diff_note={concepts: temperature dependence of the rate constant, Arrhenius equation | steps: relate the given sigma to Ea; activation energy at each temperature; sign of the temperature effect; evaluate dK/dT}]
In order to study effect of temperature on rate of reaction, another term ``$\sigma$'' is defined as $\sigma=\dfrac1K\dfrac{dK}{dT}$ [where $K$ is rate constant and $T$ represents temperature]. If for a reaction value of $\sigma=\dfrac{1.25\times10^6}{T^3}$ [$T$ should be in Kelvin], then identify the correct statement(s).
\begin{qoptions}
\opt{A}{Activation energy for the given reaction at 250 Kelvin is 10 Kcal/mol.}
\opt{B}{Activation energy for the given reaction at 2000 Kelvin is 1.25 Kcal/mol.}
\opt{C}{Rate of increase of rate constant with temperature is higher at lower temperature range than at higher temperature range.}
\opt{D}{The value of $\dfrac{dK}{dT}$ at 2000 Kelvin is $0.625\times K$.}
\end{qoptions}
\end{bankq}
\begin{bankq}[id=TNAS04-P2-CHEM-S3-Q01, source={TNAS - Test 04 (18 Jun 2023)}, paper=2, qno=1, type=numerical, subtopic={Order}, answer={15.00}, answer_src=key, difficulty=8, diff_parts={K2 S3 I1 C2 F1}, diff_note={concepts: boiling point elevation, first order kinetics | steps: moles of solute from the elevation; fraction inverted; rate constant; time for the required fraction; convert; answer}]
A 15\% by mass of sucrose solution has undergone partial inversion\par
Sucrose (aq) + Water $\longrightarrow$ Glucose(aq) + Fructose (aq)\par
After 10 hours from the start of reaction the elevation in boiling point of solution is $\dfrac{103}{200}$ K\par
The reaction follows $1^{\mathrm{st}}$ order Kinetics with respect to Sucrose. Find time (in hours) for 87.5\% of inversion of sucrose ? Neglect the change in mass of water due to reaction.\par
(Given : $K_b$(water) = 0.57 kg K mol$^{-1}$, Molar mass of sucrose = 342 gm/mol)
\end{bankq}
\begin{bankq}[id=TNAS05-P1-CHEM-S1-Q04, source={TNAS - Test 05 (09 Jul 2023)}, paper=1, qno=4, type=single, subtopic={Order}, answer={B}, answer_src=key, difficulty=4, diff_parts={K1 S2 I1 C1 F0}, diff_note={concepts: number of half lives for a given conversion | steps: fraction left; powers of two; number of half lives}]
At least how many half-lives should elapse for a $1^{\mathrm{st}}$ order reaction A $\to$ products so that the reaction is at least 95\% completed ? ($\log2=0.3$)
\begin{qoptions}
\opt{A}{4}
\opt{B}{5}
\opt{C}{6}
\opt{D}{7}
\end{qoptions}
\end{bankq}
\begin{bankq}[id=TNAS05-P1-CHEM-S2-Q02, source={TNAS - Test 05 (09 Jul 2023)}, paper=1, qno=2, type=numerical, subtopic={Rate laws}, answer={92}, answer_src=key, difficulty=8, diff_parts={K2 S3 I1 C2 F1}, diff_note={concepts: first order gas phase decomposition, vapour density | steps: fraction left after two half lives; moles of each gas; total moles; average molar mass; vapour density}]
The gas phase decomposition of dimethylether follows first order kinetics\par
$\mathrm{CH_3{-}O{-}CH_3(g)\longrightarrow CH_4(g)+H_2(g)+CO(g)}$\par
The reaction is carried out in constant volume container at $500^\circ$C and has a half-life of 14.5 min. Initially only dimethylether is present. Find $Z$ after 29 min.\par
\textbf{($Z$ = vapour density of mixture $\times$ 10)}
\end{bankq}
\begin{bankq}[id=TNAS05-P1-CHEM-S2-Q03, source={TNAS - Test 05 (09 Jul 2023)}, paper=1, qno=3, type=numerical, subtopic={Order}, answer={8}, answer_src=key, difficulty=5, diff_parts={K1 S2 I1 C1 F1}, diff_note={concepts: rate constant from initial rate data | steps: order in A; order in B; rate law constant; rate of disappearance of A}]
For the reaction $\mathrm{2A(g)+B(g)\longrightarrow C(g)+D(g)}$, at 400 K, following data are observed. Find out the rate constant (in atm$^{-2}$ sec$^{-1}$) of disappearance of A(g).\par
{\footnotesize\setlength{\tabcolsep}{3pt}\begin{tabular}{|c|p{0.17\linewidth}|p{0.17\linewidth}|p{0.17\linewidth}|}
\hline
\textbf{S.No.} & \textbf{Initial pressure of A(g)} & \textbf{Initial pressure of B(g)} & \textbf{rate of reaction (atm/sec)}\\ \hline
1. & 0.1 atm & 0.4 atm & $16\times10^{-3}$\\ \hline
2. & 0.05 atm & 0.4 atm & $4\times10^{-3}$\\ \hline
3. & 0.1 atm & 0.8 atm & $32\times10^{-3}$\\ \hline
\end{tabular}}
\end{bankq}
\begin{bankq}[id=TNAS05-P2-CHEM-S1-Q01, source={TNAS - Test 05 (09 Jul 2023)}, paper=2, qno=1, type=multi, subtopic={Order}, answer={A,C}, answer_src=key, difficulty=8, diff_parts={K2 S3 I1 C2 F1}, diff_note={concepts: first order gas phase reaction, rate of reaction | steps: order from the data; rate constant; rate at twenty minutes; rate of C; compare}]
For the gaseous reaction $A(g)\to 3B(g)+2C(g)$, following data are given.\par
{\footnotesize\begin{tabular}{|l|c|c|c|}
\hline
Time (minutes) & 10 & 20 & 60\\ \hline
Concentration of $A(g)$ & 10 M & 5 M & 2.5 M\\ \hline
\end{tabular}}\par
Select the correct statement(s).
\begin{qoptions}
\opt{A}{Rate of reaction at 20 minutes is 0.1875 M/sec}
\opt{B}{Rate of appearance of $C$ at 20 minutes is 1.5 M/sec.}
\opt{C}{Order of reaction is 3.}
\opt{D}{Order of reaction is 2.}
\end{qoptions}
\end{bankq}
\begin{bankpassage}[id=TNAS06-P2-CHEM-S3-G3-4, source={TNAS - Test 06 (30 Jul 2023)}, questions={3-4}]
An analyst starts a first order chemical reaction at 8.00 A.M. in the morning at the laboratory temperature of $27^\circ$C. At 1.00 P.M., he discovered that only 10\% of the reaction was complete by that time. To speed-up the reaction he increased the temperature to $127^\circ$C. At 4.00 P.M. he found that only 50\% of the reaction was complete from begining. Any how he did not want to stay in laboratory beyond 5.00 P.M. but he could not leave the laboratory until the reaction was 90\% complete. Fortunately he found a suitable catalyst, adding which at 4.00 P.M. at $127^\circ$C, he could meet the target of 5.00 P.M. Answer the following questions based on the above observation.\par
(Use : $\ln\dfrac{10}9=0.1$, $\ln\dfrac95=0.6$ , $\ln10=2.3$ , $\ln5=1.6$, $\ln8=2$)
\end{bankpassage}
\begin{bankq}[id=TNAS06-P2-CHEM-S3-Q03, source={TNAS - Test 06 (30 Jul 2023)}, paper=2, qno=3, type=numerical, subtopic={Arrhenius equation}, answer={5.52}, answer_src=key, difficulty=8, diff_parts={K1 S3 I2 C2 F1}, diff_note={concepts: Arrhenius equation from first order data | steps: rate constant at the low temperature; remaining fraction; rate constant at the high temperature; ratio; Arrhenius relation; activation energy}, group=TNAS06-P2-CHEM-S3-G3-4]
What was the activation energy of the original pathways (in kcal/mol) ?
\end{bankq}
\begin{bankq}[id=TNAS06-P2-CHEM-S3-Q04, source={TNAS - Test 06 (30 Jul 2023)}, paper=2, qno=4, type=numerical, subtopic={Arrhenius equation}, answer={3.92}, answer_src=key, difficulty=8, diff_parts={K1 S3 I2 C2 F1}, diff_note={concepts: catalysed activation energy from first order data | steps: rate constant after the catalyst; ratio with the uncatalysed constant; Arrhenius relation; difference of energies; subtract}, group=TNAS06-P2-CHEM-S3-G3-4]
What was the activation energy of the catalyzed pathway (in kcal/mol) ?
\end{bankq}
\begin{bankq}[id=TNAS07-P1-CHEM-S1-Q09, source={TNAS - Test 07 (20 Aug 2023)}, paper=1, qno=9, type=single, subtopic={Arrhenius equation}, answer={C}, answer_src=key, difficulty=4, diff_parts={K1 S2 I1 C1 F0}, diff_note={concepts: activation energies of forward and reverse reactions | steps: relation of enthalpy; ratio of activation energies; reverse activation energy}]
In an exothermic reaction $A(g)\rightleftharpoons B(g)$, the activation energy of reverse reaction is twice that of forward reaction. If enthalpy of reaction is $-90$ kJ/mol, the activation energy of the reverse reaction is :
\begin{qoptions}
\opt{A}{60 kJ/mol}
\opt{B}{45 kJ/mol}
\opt{C}{180 kJ/mol}
\opt{D}{90 kJ/mol}
\end{qoptions}
\end{bankq}

\section*{Star Review (Other) papers}
\begin{bankq}[id=STAR03-P2-CHEM-S1-Q03, source={Star Review (Other) - Test 03 (04 Dec 2023)}, paper=2, qno=3, type=multi, subtopic={Order}, answer={A,B,C}, answer_src=key, difficulty=6, diff_parts={K1 S2 I2 C1 F1}, diff_note={concepts: rate law and reaction order | steps: effect of each concentration; molecularity versus order; elementary steps}]
Which statement(s) is/are true for the reaction.\par
$\mathrm{C+D\longrightarrow F+G}$\par
Experimental rate law for this reaction is found to be : Rate $=K[C]^2[D]^{1/2}$
\begin{qoptions}
\opt{A}{If the concentration of $D$ is increased to four times, then the reaction rate will double.}
\opt{B}{Reaction must be a complex reaction.}
\opt{C}{Multiple elementary steps are involved in the reaction mechanism.}
\opt{D}{If the concentrations of $C$ and $D$ both are doubled, then the reaction rate will double.}
\end{qoptions}
\end{bankq}

\section*{Nurture (General) papers}
\begin{bankq}[id=NUR01-P1-CHEM-S3-Q05, source={Nurture (General) - Test 01 (12 Jun 2021)}, paper=1, qno=5, type=single, subtopic={Rate laws}, answer={7}, answer_src=solved, working={k(f)=k(b)=5x10\^{}-4 from ln(1.25)=2kt; K=1, delta H=delta G standard=0, equal activation energies, equal entropies; only E is false (and delta G of the non-equilibrium process I is not zero)}, difficulty=8, diff_parts={K3 S3 I2 C1 F0}, diff_note={concepts: racemisation kinetics and thermodynamics, rate laws, first order approach to equilibrium | steps: first order approach to equilibrium; rate constants; equilibrium constant; thermodynamic quantities; count the true statements}]
The formation of dextrorotatory isomer A from its levorotatory isomer B in water is reversible reaction having 1st order kinetics in both direction. Racemic mixture is the equilibrium product if reaction starts from either A or B. Analysis of polarimeter experiment showed that 10\% of B was converted into A in 230.3 sec at 25$^\circ$C.\par
(log 5 = 0.7, log 3 = 0.48 and log 2 = 0.3)\par
Based on given information select number of correct statement(s) among following.
\begin{qoptions}
\opt{A}{Rate constant of forward reaction at 25$^\circ$C is $5\times10^{-4}\,\mathrm{s^{-1}}$}
\opt{B}{Rate constant of backward reaction at 25$^\circ$C is $5\times10^{-4}\,\mathrm{s^{-1}}$}
\opt{C}{Activation energy of forward reaction is equal to activation energy of backward reaction.}
\opt{D}{Equilibrium constant ($K$) at 25$^\circ$C of $\mathrm{B\rightleftharpoons A}$ is 1.}
\opt{E}{Equilibrium constant increases by increasing temperature.}
\opt{F}{Molar entropy of A and B at 25$^\circ$C must be same.}
\opt{G}{$\Delta_rH$ of the process is zero.}
\opt{H}{$\Delta_rG^\circ$ of the process is zero.}
\opt{I}{$\Delta_rG$ of the process $\mathrm{B_{(l\text{-}form)}}\mathrel{\underset{k_b}{\overset{k_f}{\rightleftharpoons}}}\mathrm{A_{(d\text{-}form)}}$ is zero.}
\end{qoptions}
\end{bankq}
\begin{bankq}[id=NUR01-P2-CHEM-S1-Q07, source={Nurture (General) - Test 01 (12 Jun 2021)}, paper=2, qno=7, type=multi, subtopic={Nuclear chemistry}, answer={A}, answer_src=solved, working={27Al+4He to 30Si gives a proton, to 30P gives a neutron, 30P to 30Si gives a positron}, difficulty=3, diff_parts={K1 S1 I1 C0 F1}, diff_note={concepts: artificial transmutation | steps: conserve mass and charge; identify X, Y and Z}]
Bombardment on nucleus A by $\alpha$-particle leads to its artificial disintegration in two ways, (i) and (ii) as shown. Products X, Y and Z respectively are :\par
\bfig{126.9}{NUR01_p2_chem_s1_q07_1.png}
\begin{qoptions}
\opt{A}{proton, neutron, positron}
\opt{B}{neutron, positron, proton}
\opt{C}{proton, positron, neutron}
\opt{D}{positron, proton, neutron}
\end{qoptions}
\end{bankq}
\begin{bankq}[id=NUR03-P1-CHEM-S1-Q12, source={Nurture (General) - Test 03 (27 Sep 2023)}, paper=1, qno=12, type=single, subtopic={Rate laws}, answer={A}, answer_src=solved, working={Doubling NO gives 4 times the rate (second order), doubling H2 gives twice (first order): order 3}, difficulty=2, diff_parts={K1 S1 I1 C0 F0}, diff_note={concepts: order from initial rates | steps: compare the rates for each reactant; overall order}]
For the reaction 2NO (g) + H$_2$ (g) $\to$ N$_2$O (g) + H$_2$O (g) at 900K. following data are observed.\par
\begin{center}
{\footnotesize\begin{tabular}{|p{2.6cm}|p{2.6cm}|c|}
\hline Initial pressure of NO(atm) & Initial pressure of H$_2$(atm) & Rate\\
\hline 0.15 & 0.40 & 0.02\\
\hline 0.075 & 0.40 & 0.005\\
\hline 0.150 & 0.2 & 0.01\\
\hline
\end{tabular}}
\end{center}
Find out the rate law and order of reaction-
\begin{qoptions}
\opt{A}{k[P$_{\mathrm{NO}_2}$]$^2$[P$_{\mathrm{H}_2}$]$^1$, 3}
\opt{B}{k[P$_{\mathrm{NO}_2}$]$^2$[P$_{\mathrm{H}_2}$]$^1$, 2}
\opt{C}{k[P$_{\mathrm{NO}_2}$]$^2$[P$_{\mathrm{H}_2}$]$^1$, 1}
\opt{D}{k[P$_{\mathrm{NO}_2}$]$^2$[P$_{\mathrm{H}_2}$]$^1$, 0}
\end{qoptions}
\end{bankq}
\begin{bankq}[id=NUR03-P1-CHEM-S1-Q13, source={Nurture (General) - Test 03 (27 Sep 2023)}, paper=1, qno=13, type=single, subtopic={Rate laws}, answer={B}, answer_src=solved, working={Fe(II) and both oxalates lose 5 electrons in all, so one MnO4- reacts and 8 H+ are consumed}, difficulty=6, diff_parts={K2 S2 I2 C1 F0}, diff_note={concepts: stoichiometry of oxidation by permanganate, rate laws | steps: electrons lost by the complex; electrons gained; protons}]
In dilute aqueous H$_2$SO$_4$, the complex diaquodioxalatoferrate(II) is oxidized by MnO$_4^-$. For this reaction, the ratio of the rate of change of [H$^+$] to the rate of change of [MnO$_4^-$] is:
\begin{qoptions}
\opt{A}{16}
\opt{B}{8}
\opt{C}{4}
\opt{D}{6}
\end{qoptions}
\end{bankq}
\begin{bankq}[id=NUR03-P1-CHEM-S1-Q14, source={Nurture (General) - Test 03 (27 Sep 2023)}, paper=1, qno=14, type=single, subtopic={Order}, answer={C}, answer_src=solved, working={t50 independent of concentration: first order, t1/2=0.693/k}, difficulty=2, diff_parts={K1 S1 I1 C0 F0}, diff_note={concepts: half life against concentration plot | steps: dependence of t1/2; order}]
A graph plotted between $\log t_{50\%}$ vs. $\log$ concentration is a straight line. What conclusion can you draw from this graph, (assuming order of reaction is `n')\par
\bfig{78.1}{NUR03_p1_chem_s1_q14_1.png}
\begin{qoptions}
\opt{A}{n = 1 : t$_{1/2}\propto$ a}
\opt{B}{n = 2 : t$_{1/2}\propto$ 1/a}
\opt{C}{n = 1 : t$_{1/2}$ = (0.693/k)}
\opt{D}{None of these}
\end{qoptions}
\end{bankq}
\begin{bankq}[id=NUR03-P1-CHEM-S1-Q15, source={Nurture (General) - Test 03 (27 Sep 2023)}, paper=1, qno=15, type=single, subtopic={Rate laws}, answer={D}, answer_src=solved, working={Kc=(1e-5)\^{}2/1=1e-10; kb=kf/Kc=1.57e7}, difficulty=5, diff_parts={K2 S2 I1 C1 F0}, diff_note={concepts: rate constant of the reverse reaction, rate laws | steps: equilibrium concentrations; Kc; kb}]
The forward rate constant for the elementary reversible gaseous reaction C$_2$H$_6\rightleftharpoons$ 2CH$_3$ is $1.57\times10^{-3}$s$^{-1}$ at 100K\par
What is the rate constant for the backward reaction at this temperature if $10^{-4}$ mole of CH$_3$, 10 moles of C$_2$H$_6$ are present in a 10 litre vessel at equilibrium.
\begin{qoptions}
\opt{A}{$1.57\times10^{9}$Lmol$^{-1}$s$^{-1}$}
\opt{B}{$1.57\times10^{10}$Lmol$^{-1}$s$^{-1}$}
\opt{C}{$1.57\times10^{11}$Lmol$^{-1}$s$^{-1}$}
\opt{D}{$1.57\times10^{7}$Lmol$^{-1}$s$^{-1}$}
\end{qoptions}
\end{bankq}
\begin{bankq}[id=NUR03-P1-CHEM-S1-Q16, source={Nurture (General) - Test 03 (27 Sep 2023)}, paper=1, qno=16, type=single, subtopic={Order}, answer={B}, answer_src=solved, working={At high CA the rate tends to K1/K2: zero order}, difficulty=2, diff_parts={K1 S1 I1 C0 F0}, diff_note={concepts: order at high concentration | steps: limit of the rate law; order}]
The rate of reaction depends on concentration as $-\dfrac{dC_A}{dt}=\dfrac{K_1C_A}{1+K_2C_A}$. If C$_A$ being very high, then order of reaction is....
\begin{qoptions}
\opt{A}{1}
\opt{B}{Zero}
\opt{C}{2}
\opt{D}{3}
\end{qoptions}
\end{bankq}
\begin{bankq}[id=NUR03-P1-CHEM-S1-Q20, source={Nurture (General) - Test 03 (27 Sep 2023)}, paper=1, qno=20, type=single, subtopic={Order}, answer={C}, answer_src=solved, working={Both partial pressures change linearly with time: zero order, A giving 3B}, difficulty=6, diff_parts={K2 S2 I2 C1 F0}, diff_note={concepts: partial pressure graph, order | steps: linear change; stoichiometry; order}]
If for a reaction in which A(g) converts to B(g) the reaction carried out at const. V \& T results into the following graph.\par
\bfig{107.5}{NUR03_p1_chem_s1_q20_1.png}
\begin{qoptions}
\opt{A}{then the reaction must be A(g) $\to$3B(g) and is a first order reaction.}
\opt{B}{then the reaction must be A(g) $\to$3B(g) and is a second order reaction.}
\opt{C}{then the reaction must be A(g) $\to$3B(g) and is a zero order reaction.}
\opt{D}{then the reaction must be A(g) $\rightleftharpoons$3B(g) and is a first order reaction.}
\end{qoptions}
\end{bankq}
\begin{bankq}[id=NUR03-P1-CHEM-S2-Q01, source={Nurture (General) - Test 03 (27 Sep 2023)}, paper=1, qno=1, type=numerical, subtopic={Rate laws}, answer={4.6}, answer_src=solved, working={P=1+x/2 gives 0.9 atm reacted, 0.1 left; t=(1/5e-4)ln10=4.6x10\^{}3 s}, difficulty=7, diff_parts={K2 S2 I2 C1 F1}, diff_note={concepts: pressure of a decomposition at constant volume, rate laws | steps: pressure change with stoichiometry; first order law; time}]
The decomposition reaction\par
2N$_2$O$_5$ (g) $\xrightarrow{\Delta}$ 2N$_2$O$_4$ (g) + O$_2$ (g)\par
is started in a closed cylinder under isothermal isochoric condition at an initial pressure of 1 atm. After Y$\times10^3$ s, the pressure inside the cylinder is found to be 1.45 atm. If the rate constant of the reaction is $5\times10^{-4}$s$^{-1}$, assuming ideal gas behaviour, the value of Y is
\end{bankq}
\begin{bankq}[id=NUR03-P1-CHEM-S2-Q02, source={Nurture (General) - Test 03 (27 Sep 2023)}, paper=1, qno=2, type=numerical, subtopic={Rate laws}, answer={40}, answer_src=solved, working={k=(1/120)ln(40/5)=0.0173 per min; t1/2=40 min}, difficulty=7, diff_parts={K2 S2 I2 C1 F1}, diff_note={concepts: ester hydrolysis by titration, rate laws | steps: titre differences; first order constant; half life}]
An acid-catalysed hydrolysis of ester is carried out in the presence of HCl. Small but same volume of reaction mixture were taken out at various instant and were titrated with standard NaOH solution.\par
The following data was obtained.\par
\begin{center}
{\footnotesize\begin{tabular}{|p{3.6cm}|c|c|c|}
\hline Time (in min) & 0 & 120 & $\infty$\\
\hline Volume(mL) of NaOH solution needed for titration & 20 & 55 & 60\\
\hline
\end{tabular}}
\end{center}
The value of half-life of reaction is\underline{\hspace{1cm}}min
\end{bankq}
\begin{bankq}[id=NUR03-P1-CHEM-S2-Q09, source={Nurture (General) - Test 03 (27 Sep 2023)}, paper=1, qno=9, type=numerical, subtopic={Rate laws}, answer={4}, answer_src=solved, working={3/2 x 2.67 = 4.0 mol per min}, difficulty=3, diff_parts={K1 S1 I1 C0 F1}, diff_note={concepts: rates from stoichiometry | steps: ratio of coefficients; rate}]
The reaction that occurs in a breath analyser, a device used to determine the alcohol level in a person's blood stream is\par
2K$_2$Cr$_2$O$_7$ + 8H$_2$SO$_4$ + 3C$_2$H$_6$O $\to$ 2Cr$_2$(SO$_4$)$_3$ + 3C$_2$H$_4$O$_2$ + 2K$_2$SO$_4$ + 11H$_2$O\par
If the rate of appearance of Cr$_2$(SO$_4$)$_3$ is 2.67 mol min$^{-1}$ at a particular time, the rate of disappearance of C$_2$H$_6$O at the same time is.........mol min$^{-1}$(Nearest integer)
\end{bankq}
\begin{bankq}[id=NUR03-P1-CHEM-S2-Q10, source={Nurture (General) - Test 03 (27 Sep 2023)}, paper=1, qno=10, type=numerical, subtopic={Order}, answer={1.56}, answer_src=solved, working={k=1e-2; at 25 percent left rate=1e-2x0.25x0.25\^{}2=1.56e-4}, difficulty=4, diff_parts={K1 S2 I1 C0 F1}, diff_note={concepts: rate when 75 percent has reacted | steps: rate constant; remaining concentrations; rate}]
For the non-equilibrium process,\par
A + B $\to$ Products\par
the rate is first-order with respect to A and second order w.r.t B,\par
when [A] = 1.0M, [B] = 1.0M, rate = $1.0\times10^{-2}$ Ms$^{-1}$,\par
When 75\% each of reactant has been converted into products, rate of reaction is \underline{\hspace{0.8cm}}$\times10^{-4}$Ms$^{-1}$.
\end{bankq}
\end{document}
